%%
%% Copyright 2007-2024 Elsevier Ltd
%%
%% This file is part of the 'Elsarticle Bundle'.
%% ---------------------------------------------
%%
%% It may be distributed under the conditions of the LaTeX Project Public
%% License, either version 1.3 of this license or (at your option) any
%% later version.  The latest version of this license is in
%%    http://www.latex-project.org/lppl.txt
%% and version 1.3 or later is part of all distributions of LaTeX
%% version 1999/12/01 or later.
%%
%% The list of all files belonging to the 'Elsarticle Bundle' is
%% given in the file `manifest.txt'.
%%
%% Template article for Elsevier's document class `elsarticle'
%% with numbered style bibliographic references
%% SP 2008/03/01
%% $Id: elsarticle-template-num.tex 249 2024-04-06 10:51:24Z rishi $
%%
% =============================================================================
%  UPDATED AGAINST THE RUNS OF 2026-08-14 (20:51-21:01)
%
%  Sources used for every number below:
%     tex/tables/{gap,gap_nopen,feasibility,runtime}.tex
%     tex/tables/additional_{diesel,diesel_decomp,diesel_refuel,diesel_tw,
%                            la,la_policy,sensitivity}.tex
%     data_output/{paper_gap_stats,additional_sens_stats,additional_la_stats,
%                  additional_diesel_stats,additional_grid_stats}.csv
%     figures/*.png (2026-08-14)
%
%  COLOUR CONVENTION FOR THIS REVISION
%     \blue{...}   = number rewritten against the runs above
%     \green{...}  = my comment to you: the surrounding text no longer matches
%                    what the data does.  Delete before submission.
%     \red{...}    = your existing flags, left untouched
%
%  Everything not marked is byte-identical to your previous version.
% =============================================================================
\documentclass[preprint,12pt,3p]{elsarticle}

%% Use the option review to obtain double line spacing
%% \documentclass[authoryear,preprint,review,12pt]{elsarticle}

%% Use the options 1p,twocolumn; 3p; 3p,twocolumn; 5p; or 5p,twocolumn
%% for a journal layout:
%% \documentclass[final,1p,times]{elsarticle}
%% \documentclass[final,1p,times,twocolumn]{elsarticle}
%% \documentclass[final,3p,times]{elsarticle}
%% \documentclass[final,3p,times,twocolumn]{elsarticle}
%% \documentclass[final,5p,times]{elsarticle}
%% \documentclass[final,5p,times,twocolumn]{elsarticle}

%% For including figures, graphicx.sty has been loaded in
%% elsarticle.cls. If you prefer to use the old commands
%% please give \usepackage{epsfig}

\usepackage{amssymb}
\usepackage{amsmath}
\usepackage{comment}
\usepackage{natbib}
\usepackage{doi}
\usepackage{xcolor}
\usepackage{enumitem}
\usepackage{multirow}

\bibliographystyle{plainnat}
\usepackage{booktabs}
\usepackage{subcaption}
%% \usepackage{amsthm}
\usepackage{amsmath}

\newcommand{\red}[1]{\textcolor{red}{#1}}
\newcommand{\black}[1]{\textcolor{black}{#1}}
\newcommand{\orange}[1]{\textcolor{orange}{#1}}
\newcommand{\purple}[1]{\textcolor{purple}{#1}}
\definecolor{darkgreen}{rgb}{0.0, 0.3, 0.0}
\newcommand{\green}[1]{{\color{darkgreen}#1}}
% -- added for this revision: \blue{} marks a number rewritten against the runs.
%    (Your previous file already used \blue in the gap table but never defined
%    it, so the file did not compile as sent.)
\definecolor{darkblue}{rgb}{0.0, 0.0, 0.7}
\newcommand{\blue}[1]{{\color{darkblue}#1}}


%% The lineno packages adds line numbers. Start line numbering with
%% \begin{linenumbers}, end it with \end{linenumbers}. Or switch it on
%% for the whole article with \linenumbers.
%% \usepackage{lineno}
\begin{comment}
\journal{Energy}
\end{comment}
\begin{document}

\begin{frontmatter}

%% Title, authors and addresses

%% use the tnoteref command within \title for footnotes;
%% use the tnotetext command for theassociated footnote;
%% use the fnref command within \author or \address for footnotes;
%% use the fntext command for theassociated footnote;
%% use the corref command within \author for corresponding author footnotes;
%% use the cortext command for theassociated footnote;
%% use the ead command for the email address,
%% and the form \ead[url] for the home page:
%\title{Title\tnoteref{label1}}
%\tnotetext[label1]{Corresponding author}
%\author{Céline V. Pagnier}
%\ead{celine.pagnier@ntnu.no}
%\ead[url]{home page}
\fntext[label2]{Corresponding author: celine.pagnier@ntnu.no}
%\cortext[cor1]{Corresponding author: celine.pagnier@ntnu.no}
%% \affiliation{organization={},
%%             addressline={},
%%             city={},
%%             postcode={},
%%             state={},
%%             country={}}
%\fntext[label3]{}

\title{Joint scheduling of driver activities and charging for long distance electric truck operation under travel time uncertainty}



\author[inst1, label1]{Céline Pagnier}
\author[inst2]{Margaux Nattaf}
\author[inst2]{Marie-Laure Espinouse}
\author[inst3]{Cyrille Briand}
\author[inst1]{Steffen J.S. Bakker}


\affiliation[inst1]{organization={Department of Industrial Economics and Technology Management, Norwegian University of Science and
Technology},%Department and Organization
            addressline={Alfred Getz veg 3},
            city={Trondheim},
            postcode={NO-7491},
            %state={State One},
            country={Norway}}
\affiliation[inst2]{organization={Univ. Grenoble Alpes, CNRS, Institute of Engineering Univ. Grenoble Alpes, G-SCOP},%Department and Organization
            %addressline={Instituttveien 18, 2007 Kjeller},
            city={Grenoble},
            postcode={38000},
            %state={State Two},
            country={France}}

\affiliation[inst3]{organization={LAAS-CNRS, Université de Toulouse, CNRS, UT3},%Department and Organization
            %addressline={Instituttveien 18, 2007 Kjeller},
            city={Toulouse},
            postcode={31000},
            %state={State Two},
            country={France}}


%\begin{comment}

%% Abstract
\begin{abstract}
Battery electric trucks are entering long-haul freight operations, but various additional constraints compared to conventional diesel trucks discourage freight companies from transitioning their fleets.
As charging infrastructure expands, long-haul routes become increasingly feasible, and efficient schedules seek to synchronize charging stops with mandatory driver breaks. However, this synchronization remains fragile due to battery limitations and operational uncertainty, and becoming stranded due to lack of battery is unacceptable.
This paper studies the joint scheduling of driver activities, such as breaks and rests, and partial non-linear fast charging activities for a battery electric truck serving a fixed sequence of customers on a long-haul route under travel-time uncertainty.

We formulate a mixed-integer linear program capturing the Hours of Service regulations, including split breaks, reduced rests, extended driving allowances, and the interaction between charging and break. Around this model, we build and compare six decision architectures that differ in when decisions are committed and what information they use: a scenario-based rolling-horizon look-ahead policy, a myopic greedy policy, a robust plan, a budgeted robust plan, a two-stage stochastic plan with online duration recourse, and a perfect-hindsight oracle.

Experiments on 900 corridor instances show that policies that adapt the schedule en route perform significantly better than plans committed at departure.
Electrifying these corridors costs 8--10\,\% of route duration against a diesel truck on identical routes and identical travel-time realizations.
Finally, decision makers should direct charging-infrastructure investment towards more powerful chargers rather than a denser network: at the spacing already mandated by Regulation (EU) 2023/1804, charger power outweighs charging station density by a factor of roughly four in route duration reduction, although this shifts the burden onto grid connection capacity at each site.

\end{abstract}



%% Keywords
\begin{keyword}
%% keywords here, in the form: keyword \sep keyword
Electric trucks \sep TDSP \sep travel time uncertainty \sep rolling horizon
\end{keyword}

\end{frontmatter}

%% Add \usepackage{lineno} before \begin{document} and uncomment
%% following line to enable line numbers
%% \linenumbers





\section{Introduction} \label{intro}
Road freight accounts for roughly a quarter of transport $CO_2$ emissions in the European Union, and battery electric heavy-duty trucks (BEHDTs) are the leading decarbonization pathway for the segment. Their deployment in long-haul  service, however, is first constrained by battery capacity and charging infrastructure, but also by an operational factor that is absent for passenger electric vehicles: driver's regulations. Commercial drivers in the EU operate under Regulation (EC) No 561/2006 and Directive 2002/15/EC, which impose mandatory breaks after at most 4.5 h of driving, daily rest periods, and limits on shift driving time.
Because a mandatory 45-minute break is long enough for a substantial fast charge, synchronization of charging and drivers constraints is critical. Whether electric long-haul operation can match diesel route durations therefore hinges on how well charging events can be synchronized with legally required interruptions.

However, this synchronization is fragile. Travel times on long-haul corridors are uncertain, and a delay on a single leg propagates through the entire schedule: it consumes driving-time budget before the planned break location is reached, shifts arrivals outside customer service windows, and, because energy consumption depends on realized speed, perturbs the state of charge with which the truck reaches the next charging opportunity. A pre-determined schedule that nests charging inside breaks under average travel times may, in execution, force
the driver into an unacceptable event such as running out of battery if the energy consumption deviates from anticipated values.

This paper studies the joint scheduling of driver activities and charging for a battery electric truck on a fixed long-haul route with multiple customer stops, under stochastic travel times. Given the sequence of customers, the decisions are where and for how long to charge (with a nonlinear, piecewise-linearized charging curve and partial charging), where to take breaks and daily rests (including split breaks, reduced rests, and extended driving allowances, together with the rule that a sufficiently long charging stop can count as a break), so as to minimize route duration while respecting the daily provisions of EU HoS regulation, battery limits, and customer service windows.


Because decisions are taken sequentially while uncertainty is revealed leg by leg, the central methodological question is about the decision architecture: when should each decision be committed, and what information should it use? We organize our solution methods along exactly this axis. Three offline methods that commit a plan before departure: a box robust plan over an interval uncertainty set, a budgeted robust plan that protects only against a bounded number of simultaneous worst-case legs, and a two-stage stochastic plan whose activity structure is fixed but whose durations adapt online, and two online architectures that re-decide at every stop: a scenario-based rolling-horizon look-ahead policy and a myopic greedy policy representing typical driver behavior. All methods are evaluated in a common discrete-event simulation, and benchmarked against a perfect-hindsight oracle, so that the value of information and the value of adaptivity can be measured directly.


Our contributions are as follows.
First, we model the full weekly provisions of the EU Hours-of-Service (HoS) regulation in the context of electric trucks, incorporating partial charging along a non-linear charging curve. We formulate this as a MILP and validate it against an independent exact method for the HoS rules, and evaluate it on a real-life route used by an international trucking company.
Second, we analyze the impact of travel-time uncertainty through an extensive set of methods spanning different decision architectures, quantifying the effect of committing to a schedule upfront versus adapting en route as new information is revealed.
Third, we quantify the electrification penalty of long-haul routes by decomposing it into charging time, additional station-access time. This corroborates the finding of \citet{brockmann_break-and-charge_2025} that mandatory breaks provide sufficient charging windows most of the time, with charge/break coupling of \blue{70--76\%}, while showing that the access overheads of a fast-charging stop, which that work does not consider, absorb \blue{about half} of the resulting credit.
\green{Comment: two numbers in this sentence disagreed with
tex/tables/additional\_diesel.tex and with your own \S\,8.3.3. The coupling share $\Sigma g_i/\Sigma\tau^c_i$ is 73/76/70\,\%
by route class, so ``up to 80\,\%'' overstated the mean (individual instances do
reach 100\,\%, but the headline should be the mean). And the access overheads are
$+0.67/+1.32/+2.50$\,h against a coupling credit of $1.33/2.91/4.44$\,h, i.e.\ 45--56\,\%,
which is what \S\,8.3.3 already calls ``roughly half''. ``About a third'' was the only
place in the paper saying otherwise.}
Finally, we derive policy insights regarding charging-infrastructure investment along long-haul corridors. In particular, that charger power outweighs charger density by a factor of roughly four in its effect on route duration.


The remainder of the paper is organized as follows. Section \ref{sec:litreview} reviews related work. Section \ref{sec:problem_definition} defines the problem and states the modeling assumptions. Section \ref{sec:model} presents the deterministic MILP. Section \ref{sec:method} describes the simulation framework and the {six} decision architectures. Section \ref{sec:exp} details the experimental design, and Section \ref{sec:results} reports the computational results and disuss their implications. Section \ref{section:conclusion} concludes.



%%%%%%%%%%%%%%%%%%%%%%%%%%%%%%%%%%%%%%%%%%% -------------------------------------------------------------------------------


\section{Literature review} \label{sec:litreview}
\subsection{The truck driver scheduling problem}
Hours-of-service (HoS) regulations impose constraints on the working and driving schedule of commercial vehicle drivers. In the European Union, Regulation~(EC) No~561/2006 \citep{noauthor_regulation_2006} specifies maximum consecutive and daily driving times, mandatory break structure and minimum daily rest periods. A similar directive (2002/15/EC) limits total shift working time and restricts work during night hours. Together these rules define the truck driver scheduling problem (TDSP): given a fixed sequence of customer stops, find a feasible and cost-efficient driver schedule complying with all regulatory constraints.


The first work to model government HoS constraints in a scheduling context is \citet{xu_solving_2003}, who study a rich pickup-and-delivery problem under US rules and conjecture that finding a minimum-cost compliant schedule is NP-hard with multiple time windows. \citet{archetti_trip_2009} show that the US TDSP with single time windows is solvable in $O(n^3)$ time. For European rules, \citet{drexl_labelling_2010} formalize EU HoS as resource constraints in the elementary shortest path problem with resource constraints (ESPPRC), develop exact and heuristic labeling algorithms, and conjecture NP-completeness. \citet{goel_vehicle_2009} proposes a depth-first label-setting procedure for EU rules, which is extended in \citet{goel_truck_2010} to cover extended driving allowances and reduced rest exceptions. \citet{kok_dynamic_2010} present a dynamic programming heuristic for combined vehicle routing and driver scheduling under full EU social legislation, showing that the flexibility of complex break rules yields substantial cost reductions. \citet{prescott-gagnon_european_2010} embed EU driver rules as resources in a column-generation framework for the VRP with time windows, treating the TDSP as a pricing subproblem. \citet{rancourt_long-haul_2013} address long-haul routing under US HoS with tabu search, confirming that the choice of embedded scheduling algorithm significantly impacts solution quality. \citet{goel_hours_2014} compare EU, US, Canadian and Australian HoS regulations through optimization, finding that EU rules achieve the highest road safety at a moderate cost premium. Exact branch-and-price algorithms for the combined VRTDSP are developed by \citet{goel_exact_2017}, which is the first with proven optimality guarantees. This method is subsequently accelerated by \citet{tilk_bidirectional_2020} via bidirectional labeling. The most comprehensive EU treatment to date is \citet{garaix_label-setting_2024}, who develop a label-setting algorithm for the weekly TDSP and are the first to model the night-work rule, which limits shift working time to ten hours when night work is performed. Across this literature the driver schedule is computed once, against deterministic travel times. To our knowledge no exact EU TDSP formulation has been integrated in a sequential decision framework in which uncertain leg durations are revealed as the route is executed, which is the setting of the present paper.


\subsection{The electric vehicle routing problem}
The electric vehicle routing problem (EVRP) extends classical vehicle routing to account for the limited driving range of battery electric vehicles and the need for en-route recharging at charging stations. For a comprehensive survey we refer to \citet{erdelic_survey_2019}. The foundational formulation is the green VRP of \citet{erdogan_green_2012}, which routes alternative-fuel vehicles through refueling stations with full recharging. Exact branch-price-and-cut algorithms for the EVRPTW, covering four variants from single full recharges to multiple partial recharges, are developed by \citet{desaulniers_exact_2016}, who establish that allowing both multiple and partial recharges reduces routing cost and fleet size.

Partial recharging is studied by \citet{felipe_heuristic_2014} for a multi-technology green VRP and it is shown by \citet{keskin_partial_2016} to substantially reduce route duration . Depot-level charge scheduling over multiple days, accounting for time-dependent energy prices and battery degradation, is studied by \citet{pelletier_charge_2018}. The charging function for battery follows a concave nonlinear profile due to the constant-current/constant-voltage physics of lithium-ion cells. While most work simplify this assumption by assuming a linear charging function, it may also be linearized through the use of a piecewise-linear (PWL) approximation \citep{montoya_electric_2017}.

Uncertainty at charging stations has received growing attention: \citet{keskin_simulation-based_2021} address the EVRPTW with stochastic waiting times due to limited charger availability, showing that queuing uncertainty significantly affects routing decisions. Energy consumption uncertainty is addressed by \citet{pelletier_electric_2019} using a robust optimization framework, demonstrating that robust solutions come at only a modest cost increase. The bulk of EVRP research targets light-duty vehicles in urban or regional settings. Battery electric heavy-duty trucks (BEHDTs) differ qualitatively: much larger energy demands per leg, reliance on high-power fast charging, and binding HoS constraints that govern both when the driver may be on duty and how charging time interacts with mandatory break requirements.


\subsection{Combining TDSP and EVRP}
Despite the practical importance of long-haul BET operations, the literature combining HoS regulation with battery energy management in a single optimization model is sparse. \citet{bai_rollout-based_2023} study the problem of jointly determining optimal charging and rest decisions for an electric truck on a fixed pre-planned route under HoS constraints. They formulate a bilinear MIP minimizing total cost (charging plus rest time) and develop a rollout-based approximate solver that scales linearly with the number of stations, validated on Swedish road network data.

\citet{brockmann_break-and-charge_2025} develop a mathematical programming model that synchronizes driver break times with truck charging events under EU HoS, validated on real-world data from northwest Germany. Their key finding is that mandatory breaks create natural charging windows sufficient to match diesel route durations. However they do not consider additional manever time due to access to the charging stations. \citet{wan_long-haul_2024} study long-haul truck charging planning with time and energy flexibility, proposing a mixed-integer model for opportunity charging at routine stops. \citet{su_optimizing_2025} introduce a stage-expanded graph for minimizing the carbon footprint of long-haul heavy-duty e-truck trips, simultaneously optimizing path, speed, and charging, and show that the optimized charging strategy compounds the emissions benefit of electrification. \citet{vidotte_plaza_charger_2026} are the first to jointly optimize charging infrastructure siting and HoS-compliant driver scheduling for long-haul eHDT fleets, though their focus is on strategic network design rather than operational online scheduling. Finally, \citet{niyomphon_stochastic_2026} propose a chance-constrained stochastic MILP for electric freight on fixed long-haul routes with partial recharging and uncertain freight demand, but without HoS constraints.


\subsection{Uncertainty in long-haul freight transport}
Travel time uncertainty in long-haul road freight creates two risks: a HoS violation when accumulated driving or working time exceeds its limit before a scheduled break, and battery depletion when a leg uses more energy than planned. While stochastic EVRP models have addressed energy consumption uncertainty \citep{pelletier_electric_2019} and queuing uncertainty at chargers \citep{keskin_simulation-based_2021}, travel time uncertainty and its joint effect on HoS compliance and SOC feasibility remains largely unexplored.

The framework of this type of problem where decisions have to be made sequentially, and adjusted based on realized uncertainty, fits well with the rolling horizon methodology, especially when the multi stage optimization problem would be too complex to solve. The rolling-horizon paradigm is a well-established technique for multi-period optimization under gradual information revelation \citep{glomb_rolling-horizon_2022}. The broader literature on dynamic vehicle routing \citep{pillac_review_2013} covers stochastic programming, robust optimization, and online policies, and classifies our problem as one with a fully known route structure but sequentially revealed stochastic leg costs.


\subsection{Contribution}

Table~\ref{tab:litreview} positions the present work against the most closely related papers across eight dimensions. This paper is the first to jointly model full EU HoS (including split-break sequencing, shift working-time limits, multiple rest types, and charging-break interactions), partial fast-charging with a nonlinear PWL charging curve, multiple customer stops with soft time windows on a fixed long-haul route, and stochastic travel times.

{\color{red} CHECK AGAIN TABLE}

\begin{table}[htbp]
\centering
\caption{Positioning of the present work in the literature.}
\label{tab:litreview}
\scriptsize
\setlength{\tabcolsep}{3.5pt}
\resizebox{\textwidth}{!}{%
\begin{tabular}{p{3.8cm} c p{1.8cm} p{1.5cm} c c c}
\toprule
\textbf{Reference}
  & \textbf{HoS}
  & \textbf{Charging}
  & \textbf{Route length}
  & \textbf{Fixed}
  & \textbf{TW}
  & \textbf{Uncertainty}\\
\midrule

\citealt{goel_exact_2017}   & EU & $-$ & Long-haul & $-$ & Hard & $-$\\
\citealt{garaix_label-setting_2024}   & EU & $-$ & Long-haul & $\checkmark$ & Hard & $-$\\

\citealt{montoya_electric_2017}   & $-$ & Partial + NL & Regional & $-$ & Hard & $-$\\

\citealt{keskin_simulation-based_2021}   & $-$ & Partial + L & Long-haul & $-$ & Hard & CS queues\\
\citealt{niyomphon_stochastic_2026}   & $-$ & Partial + L & Long-haul & $-$ & $-$ & Demand \\

\citealt{bai_rollout-based_2023}   & EU & Partial$+$L & Day & $\checkmark$ & $-$ & $-$\\
\citealt{brockmann_break-and-charge_2025}   & EU & Partial$+$NL  & Day & $\checkmark$ & $-$ & $-$ \\
\citealt{vidotte_plaza_charger_2026}   & Brazil & Full$+$L &  Long-haul & $-$ & Hard & $-$ \\
\midrule
This paper
  & {EU}
  & {Partial$+$NL}
  & {Long-haul}
  & $\checkmark$
  & {Soft}
  & Travel time\\
\bottomrule
\end{tabular}}
\par\smallskip
\begin{minipage}{\textwidth}
\footnotesize
\textbf{Charging}: L/NL: linear/non-linear charging function.
\textbf{TW}: customer time windows.
\end{minipage}
\end{table}


\section{Problem Definition}
\label{sec:problem_definition}

This paper studies the problem of scheduling driver activities for long-distance freight transport with a battery electric truck (BET). The schedule must simultaneously satisfy Hours-of-Service (HoS) regulations and battery energy constraints. Given a predefined sequence of customer visits, the objective is to minimize the total route duration by determining the optimal timing and locations of mandatory breaks, daily rest periods, and battery charging operations.

\subsection{Route and Customers}
The route connects a sequence of customer locations via road segments whose lengths are known in advance. Travel on each segment is characterized by an average speed. In practice, however, traffic conditions introduce uncertainty in the actual travel duration: this is explicitly accounted for in the model. The total route duration is the sum of five components: driving time accumulated while the vehicle is moving between locations, service time for loading or unloading operations, charging time at charging stations, station access and maneuvering time, and legally required interruptions such as mandatory breaks and daily rest periods.

The sequence of customers to be visited is fixed and known prior to departure. At each customer location, the driver performs a loading or unloading operation whose duration is known.
At each customer, soft time windows are considered. Service begins immediately upon arrival, whether or not the window is open: distribution terminals do not turn a truck away, but an arrival outside the announced window, either early or late, disrupts terminal planning and incurs a fixed service-level penalty, identical for early and late arrivals. Any break or rest
taken at a customer stop occurs after service. The vehicle is not permitted to wait idly for a window to open: an early arrival is served early and penalized, or the driver stretches a break at an earlier stop.

\subsection{Charging operations}
The BET has a battery of fixed capacity (kWh). Energy consumption is assumed to be dependent of the distance traveled and the vehicle speed, yielding a specific consumption rate expressed in kWh/km. At each CS stop, the driver may charge the battery either partially or fully, and multiple charging stops are permitted throughout the route. Charging is not allowed during customer service operations.

Only high-power fast charging is available along the route (overnight slow charging such as in depot is not considered). The charging process is modeled by a piecewise linear approximation of the nonlinear charge-rate curve. Queuing at a CS is modeled as a known waiting time at each CS.

Finally, a sufficiently long charging stop may simultaneously satisfy a mandatory break requirement. Specifically, a charging event counts as a driving break provided its duration meets the applicable threshold: at least 15~minutes for a first split break, at least 30~minutes for a second split break, or at least 45~minutes for a full break, as defined by the HoS regulations described below.

\subsection{HoS Regulations}
The scheduling problem is subject to European HoS regulations given by Regulation (EC) No 561/2006 (driving times, breaks and rests) and Directive 2002/15/EC (working time).
Two accumulation counters are maintained: driving time, which counts only periods when the vehicle is in motion, and working time, which additionally includes customer service, charger access and maneuvering, and charging time not covered by a simultaneous break.

Regarding mandatory breaks, a driver may not accumulate more than 4.5~hours of continuous driving without taking a full break of at least 45~minutes. This full break may be replaced by a split sequence consisting of a first period of at least 15~minutes (but less than 45~minutes) followed, in that order, by a second period of at least 30~minutes.

Regarding shift-level constraints, the maximum permitted driving time per shift is 9~hours, extendable to 10~hours on no more than two occasions per week.
%Furthermore, if the shift includes a period of night working, the total working time within that shift is limited to 10~hours.
A daily rest of at least 11 h must separate consecutive shifts, and it may be reduced to 9 h at most three times between weekly rests.
As a consequence of these rest requirements, the sum of working time, break time, and any other on-duty periods within a 24-hour window is bounded by 13~hours under a full rest, or by 15~hours when the rest is reduced to 9~hours.

Charging during rests is treated as sequential: the truck must vacate the charging bay before the rest begins, and such charging time counts as work.

While charging is only possible at charging stations,  breaks and rests may be taken at customer stops, charging stations, and dedicated rest-area nodes inserted along legs.


\section{Deterministic Mathematical Model} \label{sec:model}
We now present the mathematical formulation for the deterministic problem. Tables ~\ref{tab:notation-mini} and ~\ref{tab:notation-mini2} present all notation used in the model formulation. Figure \ref{fig:solution} present a possible solution of the problem.



\begin{table}[h!]
\centering
\caption{Notation summary.}
\label{tab:notation-mini}
\small
\begin{tabular}{p{0.12\linewidth}p{0.8\linewidth}}
\toprule
\multicolumn{2}{l}{\textbf{Sets}}\\
\midrule
$\mathcal{I}$ & Set of possible stops, ordered from 0=origin to N=destination.\\
$\mathcal{C}$ & Set of customers $\mathcal C \subseteq \mathcal I$.\\
$\mathcal{K}$ & Set of charging stations $\mathcal K \subseteq \mathcal I$.\\
$\mathcal{L}$ & Set of rest areas $\mathcal L \subseteq \mathcal I$, $\mathcal C \cap \mathcal K \cap \mathcal L= \emptyset$.\\
$\mathcal{R}$ & Set of breakpoints of the piecewise-linear charging function.\\
\midrule
\multicolumn{2}{l}{\textbf{Parameters}}\\
\midrule
$D_i$ & Travel time from stop $i$ to stop $i+1$.\\
$E_i$ & Energy consumption between stop $i$ and stop $i+1$.\\
$E^{cap}$ & Battery capacity. \\
$E^{\min}$ & Minimum SOC. \\
$E^0$ & Initial SOC. \\

$\bar{E}_r$  & SOC level at breakpoint $r \in \mathcal{R}$ of the charging fct, ($\bar{E}_0 = 0$, $\bar{E}_R = E^{cap}$).\\
$\mathcal{T}_r$ & Time to charge from empty to $\bar{E}_r$,($\mathcal{T}_0 = 0$).\\
$\mathcal{T}_{\text{min}}$ & Minimum charging time at a CS.\\


$S_i$ & Service time at customer $i \in \mathcal{C}$. \\
$Q_i$ & Queue time at CS stop $i \in \mathcal{K}$ if charging occurs. Zero if $y_i = 0$.\\
$M_i^{stop}$ & Maneuver time at CS stop $i \in \mathcal{S}$ if any activity occurs.\\
$M_i^{seq}$ & Repositioning time at stop $i \in \mathcal{K}$ when charging and a break or rest are declared sequentially, i.e. the truck must vacate the charging bay before the break begins.\\

$W_i^{Ha}, W_i^{Hf}$  & Service time window at customer $i \in \mathcal{C}$. \\
$ \beta$ & Fixed out-of-window service penalty (h). \\

%$[W_i^{Sa}, W_i^{Sf}]$  & Soft time window at customer $i \in \mathcal{C}$ \\
$T^{brk}_{b45}$   & Full break duration ($= 45$\,min). \\
$T^{brk}_{b_{15}}$     & Min first split break ($= 15$\,min). \\
$T^{brk}_{b_{30}}$     & Min second split break ($= 30$\,min). \\
$T^{rst}_{r_1}$          & Minimum daily rest duration ($= 11$\,h). \\
$T^{rst}_{r_2}$          & Reduced daily rest duration (exception 1) ($= 9$\,h at most $3\times$/week). \\
$\bar{\rho}$ & Weekly number of allowed reduced rest duration ($=3$).\\

$T^{drv}_{cons}$   & Max consecutive driving time ($= 4.5$\,h). \\
$T^{drv}_{shift_1}$     & Max shift driving time ($= 9$\,h). \\
$T^{drv}_{shift_2}$     & Increased shift driving time (exception 2) ($= 10$\,h at most $2\times$/week). \\
$\bar \epsilon $ &  Weekly number of allowed extended driving duration ($=2$).\\

$T^{wk}_{sh1}$      & Max shift working time under a regular daily rest ($= 13$\,h). \\
$T^{wk}_{sh2}$      & Max shift working time under a reduced daily rest ($= 15$\,h). \\

\bottomrule
\multicolumn{2}{l}{{\footnotesize Abbreviations: SOC = state of charge, CS = charging station}}\\
\end{tabular}
\end{table}


\begin{table}[h!]
\centering
\caption{Variables notation summary.}
\label{tab:notation-mini2}
\small
\begin{tabular}{p{0.25\linewidth}p{0.75\linewidth}}
\toprule
\multicolumn{2}{l}{\textbf{Variables}}\\
\midrule
$t_i^a,t_i^d \geq 0$ & arrival/departure time at stop $ i \in \mathcal{I}$. \\
$\delta_i \in \{0,1\}$ & 1 if service at customer $ i \in \mathcal{C}$ starts outside its window.\\

$y_i \in \{0,1\}$ & 1 if the truck charges at stop $ i \in \mathcal{K}$.\\
$\tau_i^c \geq 0$ & charging duration at stop $ i \in \mathcal{K}$.\\
$e_i^a, e_i^d \geq 0$ & SOC at arrival/departure at CS at stop $ i \in \mathcal{I}$.\\

$\lambda_{i,r}^{a},\lambda_{i,r}^{d} \geq 0$ & PWL weight for encoding $e_i^a$, $e_i^d$, $i \in \mathcal{K}$, $r \in \mathcal{R}$.\\
$\mu_{i,r}^{a},\mu_{i,r}^{d} \in \{0,1\}$ & $1$ if $e_i^a$,$e_i^d$ lies in segment $r$ of the PWL curve, $i \in \mathcal{K}$, $r \in \mathcal{R}$.\\

$c_i^d \geq 0$ & accumulated consecutive driving time at arrival at stop $i$ (h).\\
$s_i^d \geq 0$ & accumulated shift driving time at arrival at stop $i$ (h).\\
$s_i^w \geq 0$ & accumulated shift working time at arrival at stop $i$ (h).\\
$h_i \geq 0$ & Elapsed time since end of last daily rest (h).\\

$x_i^{b_{45}}, x_i^{b_{15}}, x_i^{b_{30}} \in \{0,1\}$ & break type indicators at stop $i$. \\
$\phi_i \in \{0,1\}$&  1 if the first split break has been taken since the last driving reset.\\
$\tau_i^{b} \geq 0$& Break time beyond charging at stop $i$. \\
$\hat{\tau}_i^{b} \geq 0$& Break duration credited for HoS purposes. \\
$g_i \geq 0$ & Charging time credited as break at stop $i$. \\
$\sigma_i \in \{0,1\}$ & 1 if CS $i$ operates in sequential mode (no charging/break synchronization).\\


$\rho_i^{r_1}, \rho_i^{r_2} \in \{0,1\}$& rest type indicators at stop $i$.\\
$\tau_i^r \geq 0$& rest duration at stop $i$.\\

$v_i \in \{0,1\}$ & 1 if any activity (charge/break/rest) occurs at stop $i$.\\
$z_i, q_i \in \{0,1\}$ & extended shift flag; extension consumed at a rest.\\

\bottomrule
\multicolumn{2}{l}{{\footnotesize Abbreviations: HoS = Hours-of-service, CS = Charging station}}\\
\end{tabular}
\end{table}



\subsection{Sets and Parameters}
Let the route consist of a fixed sequence of stops indexed by $i \in \mathcal{I} = \{0, 1, \ldots, N\}$, where stop $0$ is the origin and stop $N$ is the destination. The set of customer stops is $\mathcal{C} \subseteq \mathcal{I}$, and the set of CS stops is $\mathcal{K} \subseteq \mathcal{I}$. Additionally, rest areas $\mathcal{L} \subseteq \mathcal{I}$ are place along the route and are treated as Charging station but with charging disabled. The set of non-customer stops $\mathcal S$ include both charging station and rest areas, such that  $\mathcal{C} \cap \mathcal{K} \cap \mathcal{L} = \emptyset$ and $\mathcal{K} \cup \mathcal{L} = \mathcal S$.  The model structure is similar to that of \citet{bai_rollout-based_2023}, adapted here to a single-vehicle fixed-route setting with multiple customer stops, as seen in Figure \ref{fig:routeExample}.
The nonlinear charging curve is approximated by a piecewise linear function defined by $R+1$ breakpoints, following \citet{montoya_electric_2017}.


\begin{figure}
    \centering
    \includegraphics[width=1\linewidth]{figures/routeExample.png}
    \caption{Structure of a benchmark corridor. Charging stations (triangles) are placed on a 60 km grid,  rest-area nodes (ticks) subdivide any leg longer than 25 km and allow for a break or rest but no charging, customers (diamonds) are drawn around cluster centers at 25\%, 50\% and 70\% of the corridor. The instance shown is a short-route case: 900 km, three customers.}
    \label{fig:routeExample}
\end{figure}




\subsection{Objective Function}
The objective minimizes the arrival time at the destination plus a fixed penalty $\beta$ for every customer whose service starts outside its window.
\begin{equation}
    \min \quad t_N^a + \beta \sum_{i \in \mathcal C}\delta_i
\end{equation}

\subsection{Constraints}
In order to simplify the following constraints, we define additional variables as such: $\rho_i = \rho_i^{r1} + \rho_i^{r2} \in \{0,1\}$ is the rest indicator at a stop, $x_i = x_i^{b45} + x_i^{b15} + x_i^{b30} \in \{0,1\}$ is the break indicator at a stop, $r_i = x_i^{b45} + x_i^{b30} + \rho_i^{r1} + \rho_i^{r2} \in \{0,1\}$ is the reset indicator of consecutive drive time at a stop. We further write $o_i=t^d_i - t^a_i - \tau_i^r $ for the on-duty (non-rest) time spent at stop $i$, used in the worktime shift counter.

\subsubsection{Time propagation}
Constraint~\eqref{eq:time_prop} states that the arrival time at each stop equals the departure time from the previous stop plus the leg travel time. The departure time from a stop depends on the activities undertaken there. At a customer stop, constraint~\eqref{eq:depart_C} accounts for mandatory service time $S_i$ and any additional break or rest duration. At a CS stop, constraint~\eqref{eq:depart_K} additionally includes a queue time $Q_i y_i$, a charging duration $\tau_i^c$, a maneuver time $M_i^{stop}$ to get to/leave the CS in case we stop there for any activity, a maneuver time $M_i^{seq}$ if we first charge and then sequentially take a break or rest at the CS. Constraints~\eqref{eq:tw_soft1}--\eqref{eq:tw_soft2} activate the out-of-window indicator $\delta_i$ whenever service starts before the window opens or after it closes. Since service starts at arrival, the window applies to $t_i^a$.
\begin{equation}
    t_{i+1}^a = t_i^d + D_i,
    \qquad \, i \in \mathcal{I}\setminus\{N\}
    \label{eq:time_prop}
\end{equation}

\begin{equation}
    t_i^d = t_i^a + S_i + \tau_i^{b} + \tau_i^r,
    \qquad \, i \in \mathcal{C}
    \label{eq:depart_C}
\end{equation}

\begin{equation}
    t_i^d = t_i^a +  M_i^{stop} v_i + Q_i y_i + \tau_i^c + M_i^{seq} \sigma_i + \tau_i^{b} + \tau_i^r,
    \qquad \, i \in \mathcal{S}
    \label{eq:depart_K}
\end{equation}


\begin{equation}
    t_i^a \leq W_i^{f} + M \delta_i,
    \qquad \, i \in \mathcal{C}
    \label{eq:tw_soft1}
\end{equation}
\begin{equation}
    t_i^a \geq W_i^{a} - M \delta_i,
    \qquad \, i \in \mathcal{C}
    \label{eq:tw_soft2}
\end{equation}


\subsubsection{Battery state-of-charge}
Constraint~\eqref{eq:soc_prop} propagates the SOC between stops according to the energy consumed on each leg. Constraints~\eqref{eq:soc_nc}--\eqref{eq:soc_mono} enforce that charging can only occur at designated CS stops and can never decrease the SOC. Constraints~\eqref{eq:soc_lb}--\eqref{eq:soc_ub} keep the SOC within the feasible battery operating range at all times. Constraint~\eqref{eq:chg_act} ensures that a positive charging duration can only be incurred when the binary $y_i$ is active, and that in case of charging, we charge at least as much as the minimum charging time allowed.

\begin{equation}
    e_{i+1}^a = e_i^d - E_i,
    \qquad \, i \in \mathcal{I}\setminus\{N\}
    \label{eq:soc_prop}
\end{equation}
\begin{alignat}{2}
    e_i^d &= e_i^a, &\qquad& \, i \in \mathcal{C} \cup \mathcal L \label{eq:soc_nc}\\
    e_i^d &\geq e_i^a, &\qquad& \, i \in \mathcal{K} \label{eq:soc_mono}
\end{alignat}
\begin{alignat}{2}
    e_i^a &\geq E^{\min}, &\qquad& \, i \in \mathcal{I}
    \label{eq:soc_lb}\\
    e_i^d &\leq E^{\mathrm{cap}}, &\qquad& \, i \in \mathcal{I}
    \label{eq:soc_ub}
\end{alignat}
\begin{equation}
    \mathcal{T}_{\text{min}} y_i \leq \tau_i^c \leq \mathcal{T}_R\, y_i,
    \qquad \, i \in \mathcal{K}
    \label{eq:chg_act}
\end{equation}




\subsubsection{Piecewise-Linear Charging Function}
Following \citet{montoya_electric_2017}, we linearize the nonlinear charging function using a convex-combination formulation. Constraints~\eqref{eq:pwl_ea}--\eqref{eq:pwl_ed} express both the arrival and departure SOC at each CS stop as convex combinations of the $R+1$ PWL breakpoints. The charging duration is obtained in constraint~\eqref{eq:pwl_tauc} as the difference in cumulative charge time between the departure and arrival SOC levels. Constraints~\eqref{eq:pwl_ca}--\eqref{eq:pwl_cd} impose the convex-combination normalization condition on the arrival and departure weights. Constraints~\eqref{eq:sos2_sa}--\eqref{eq:sos2_sd} enforce that exactly one PWL segment is active for each SOC level. Finally, the SOS2 adjacency constraints~\eqref{eq:sos2_lo}--\eqref{eq:sos2_hi} restrict non-zero weights to at most two consecutive breakpoints, for both arrival ($\nu=a$) and departure ($\nu=d$) SOC levels.

\begin{align}
    e_i^a &= \sum_{r \in \mathcal{R}} \lambda_{i,r}^{a}\,\bar{E}_r,
    \qquad \, i \in \mathcal{K} \label{eq:pwl_ea}\\
    e_i^d &= \sum_{r \in \mathcal{R}} \lambda_{i,r}^{d}\,\bar{E}_r,
    \qquad \, i \in \mathcal{K} \label{eq:pwl_ed}
\end{align}

\begin{equation}
    \tau_i^c = \sum_{r \in \mathcal{R}} \lambda_{i,r}^{d}\,\mathcal{T}_r
             - \sum_{r \in \mathcal{R}} \lambda_{i,r}^{a}\,\mathcal{T}_r,
    \qquad \, i \in \mathcal{K}
    \label{eq:pwl_tauc}
\end{equation}

\begin{alignat}{2}
    \sum_{r \in \mathcal{R}} \lambda_{i,r}^{a} &= 1, &\qquad& \, i \in \mathcal{K}
    \label{eq:pwl_ca}\\
    \sum_{r \in \mathcal{R}} \lambda_{i,r}^{d} &= 1, &\qquad& \, i \in \mathcal{K}
    \label{eq:pwl_cd}
\end{alignat}

\begin{alignat}{2}
    \sum_{r=1}^{R} \mu_{i,r}^{a} &= 1, &\qquad& \, i \in \mathcal{K}
    \label{eq:sos2_sa}\\
    \sum_{r=1}^{R} \mu_{i,r}^{d} &= 1, &\qquad& \, i \in \mathcal{K}
    \label{eq:sos2_sd}
\end{alignat}

\begin{align}
    \lambda_{i,0}^{\nu}   &\leq \mu_{i,1}^{\nu},
    \label{eq:sos2_lo}\\
    \lambda_{i,r}^{\nu}   &\leq \mu_{i,r}^{\nu} + \mu_{i,r+1}^{\nu},
    \quad r = 1,\ldots,R-1,
    \label{eq:sos2_mid}\\
    \lambda_{i,R}^{\nu}   &\leq \mu_{i,R}^{\nu},
    \label{eq:sos2_hi}
\end{align}
for all $i\in\mathcal{K}$ and $\nu\in\{a,d\}$.


\subsubsection{Break and Rest Decisions}
Constraint~\eqref{eq:one_act} ensures that at most one type of break or rest activity is declared at each stop.
\begin{equation}
    x_i    + \rho_i \leq 1,
    \qquad \, i \in \mathcal{I}
    \label{eq:one_act}
\end{equation}

At a CS stop, charging time may qualify toward HoS break requirements only when a break is declared and runs in parallel with charging ($x_i = 1$, $\sigma_i = 0$).
Let $g_i$ denote the credited charging time. Then constraints \eqref{eq:pi1}--\eqref{eq:pi3} force $g_i = \tau_i^c$ exactly in that case and $g_i = 0$ otherwise.
The effective break duration $\hat \tau_i^b$ is then the credited charging time plus any additional break time \eqref{eq:bhat_K}, and the additional break time alone at customer and rest-area stops \eqref{eq:bhat_C}.

\begin{equation}
    g_i \leq \tau_i^c
    \qquad \, i \in \mathcal{K}
    \label{eq:pi1}
\end{equation}
\begin{equation}
    g_i \leq \mathcal T_R (1- \sigma_i)
    \qquad \, i \in \mathcal{K}
    \label{eq:pi2}
\end{equation}
\begin{equation}
    g_i \leq \mathcal T_R x_i
    \qquad \, i \in \mathcal{K}
    \label{eq:pi21}
\end{equation}
\begin{equation}
    g_i \geq \tau_i^c - \mathcal T_R \sigma_i - \mathcal T_R (1-x_i)
    \qquad \, i \in \mathcal{K}
    \label{eq:pi3}
\end{equation}

\begin{alignat}{2}
    \hat{\tau}_i^b &= g_i+ \tau_i^b, &\quad& \, i \in \mathcal{K}
    \label{eq:bhat_K}\\
    \hat{\tau}_i^b &= \tau_i^b,            &\quad& \, i \in \mathcal{C} \cup \mathcal L
    \label{eq:bhat_C}
\end{alignat}


Constraints~\eqref{eq:brk45}--\eqref{eq:brk30} impose the minimum effective break duration for each break type. Similarly, constraints~\eqref{eq:rst1}--\eqref{eq:rst2} impose the minimum rest duration for each rest type. Constraints~\eqref{eq:brk_ub}--\eqref{eq:rst_ub} force the break and rest durations to zero when no corresponding activity is declared. Constraint~\eqref{eq:r2_lim} enforces that the reduced daily rest of $9$\,h may be used at most $\bar \rho$ times per week.

\begin{alignat}{2}
    \hat{\tau}_i^b &\geq T^{\mathrm{brk}}_{45}\, x_i^{b_{45}}, &\quad& \, i \in \mathcal{I}
    \label{eq:brk45}\\
    \hat{\tau}_i^b &\geq T^{\mathrm{brk}}_{15}\, x_i^{b_{15}}, &\quad& \, i \in \mathcal{I}
    \label{eq:brk15}\\
    \hat{\tau}_i^b &\geq T^{\mathrm{brk}}_{30}\, x_i^{b_{30}}, &\quad& \, i \in \mathcal{I}
    \label{eq:brk30}
\end{alignat}
\begin{alignat}{2}
    \tau_i^r &\geq T^{\mathrm{rst}}_{1}\, \rho_i^{r_1}, &\quad& \, i \in \mathcal{I}
    \label{eq:rst1}\\
    \tau_i^r &\geq T^{\mathrm{rst}}_{2}\, \rho_i^{r_2}, &\quad& \, i \in \mathcal{I}
    \label{eq:rst2}
\end{alignat}
\begin{alignat}{2}
    \tau_i^b &\leq M\, x_i,    &\quad& \, i \in \mathcal{I}
    \label{eq:brk_ub}\\
    \tau_i^r &\leq M\, \rho_i, &\quad& \, i \in \mathcal{I}
    \label{eq:rst_ub}
\end{alignat}
\begin{equation}
    \sum_{i \in \mathcal{I}} \rho_i^{r_2} \leq \bar{\rho}
    \label{eq:r2_lim}
\end{equation}

EU regulation allows the standard $45$\,min break to be replaced by a split break taken in two parts: a first part of at least $15$\,min followed by a second part of at least $30$\,min, in that strict order. The binary $\phi_i$ tracks whether the first part of a split break has been completed since the last consecutive-driving reset. Constraint~\eqref{eq:split_ord} ensures the second part can only be declared if the first part was taken at an earlier stop. Constraints~\eqref{eq:phi1}--\eqref{eq:phi3} propagate $\phi$ between stops. It is set to $1$ when a first split-break part is taken and no reset event has occurred, it is cleared to $0$ by any consecutive-driving reset (a full break, the completion of the second split-break part, or any rest).
\begin{equation}
    x_i^{b_{30}} \leq \phi_i,
    \qquad \, i \in \mathcal{I}
    \label{eq:split_ord}
\end{equation}

\begin{alignat}{3}
    \phi_{i+1} &\geq \phi_i + x_i^{b_{15}} - r_i,
    &\quad& \, i \in \mathcal{I}\setminus\{N\}
    \label{eq:phi1}\\
    \phi_{i+1} &\leq \phi_i + x_i^{b_{15}},
    &\quad& \, i \in \mathcal{I}\setminus\{N\}
    \label{eq:phi2}\\
    \phi_{i+1} &\leq 1 - r_i,
    &\quad& \, i \in \mathcal{I}\setminus\{N\}
    \label{eq:phi3}
\end{alignat}

Constraint~\eqref{eq:anyactivity} flags any activity at a stop of $\mathcal S$. Sequential mode $\sigma_i$ is possible only when charging occurs \eqref{eq:sigma1}. It is forced whenever charging coexists with a rest \eqref{eq:sigma2}, since the truck must vacate the charging bay before an 11\,h or 9\,h rest, and it is forced whenever a break is declared that the charge does not fully cover \eqref{eq:sigma3}, since the driver must then take the break separately after charging. In \eqref{eq:sigma3}, $\underline b_i = T^{brk}_{45}x^{b_{45}}_i + T^{brk}_{15}x^{b_{15}}_i + T^{brk}_{30}x^{b_{30}}_i$ is the required minimum of the break declared at stop $i$; if the credited charge $g_i$ (equivalently, since $g_i = \tau^c_i$ under parallel mode, the charge itself) falls short of it, $\sigma_i$ is forced to 1, and constraint~\eqref{eq:pi2} then sets $g_i = 0$ so that the full break sits sequentially after the charge.
\begin{alignat}{2}
    \tfrac{1}{2}(y_i + \rho_i + x_i) \;\leq\; v_i &\;\leq\; y_i + \rho_i + x_i,
    &\qquad& i \in \mathcal{S}
    \label{eq:anyactivity}\\
    \sigma_i &\leq y_i, &\qquad& i \in \mathcal{K}
    \label{eq:sigma1}\\
    \sigma_i &\geq y_i + \rho_i - 1, &\qquad& i \in \mathcal{K}
    \label{eq:sigma2}\\
    \sigma_i &\geq \frac{\underline b_i - \tau^c_i}{T^{brk}_{45}} - (1 - y_i), &\qquad& i \in \mathcal{K}
    \label{eq:sigma3}
\end{alignat}

% ----------------------------------------------------------



\subsubsection{Driving regulation}

The accumulated consecutive driving time $c_i^d$ increases with each driven leg and is reset to zero by any qualifying break or rest (i.e.\ whenever $r_i = 1$). Defining the auxiliary $l_i^1 = r_i\, c_i^d$ as the portion of $c_i^d$
that is zeroed out at stop $i$, constraint~\eqref{eq:cd_prop} expresses the propagation rule. Constraints~\eqref{eq:l1_ub1}--\eqref{eq:l1_lb} linearize the bilinear product $l_i^1 = r_i\, c_i^d$ via a standard big-$M$ reformulation. Constraint~\eqref{eq:cd_ub} enforces the regulatory upper bound on consecutive driving time.

\begin{equation}
    c_{i+1}^d = c_i^d + D_i - l_i^1,
    \qquad \, i \in \mathcal{I}\setminus\{N\}
    \label{eq:cd_prop}
\end{equation}

\begin{equation}
    c_i^d \leq T^{\mathrm{drv}}_{\mathrm{cons}},
    \qquad \, i \in \mathcal{I}
    \label{eq:cd_ub}
\end{equation}

The accumulated shift driving time $s_i^d$ is reset only by a daily rest (a break alone is insufficient). Defining $l_i^2 = \rho_i\, s_i^d$, constraint~\eqref{eq:sd_prop} gives the propagation rule. Constraints~\eqref{eq:l2_ub1}--\eqref{eq:l2_lb} linearize $l_i^2 = \rho_i\, s_i^d$ with big-M $T^{drv}_{sh_2} = 10$\,h, the largest value $s^d_i$ can legitimately attain on an extended shift. Constraint~\eqref{eq:sd_ub} imposes the maximum shift driving time, conditional on the extension flag $z_i$. Constraint~\eqref{eq:z_prop} lets the flag persist within a shift and reset only after a rest. Constraints~\eqref{eq:q_def}--\eqref{eq:ext_budget} charge one unit of the weekly budget $\bar\epsilon$ per extended shift, the term $z_N$ covering an unfinished final shift.
\begin{equation}
    s_{i+1}^d = s_i^d + D_i - l_i^2,
    \qquad \, i \in \mathcal{I}\setminus\{N\}
    \label{eq:sd_prop}
\end{equation}

\begin{equation}
    s_i^d \leq T^{\mathrm{drv}}_{\mathrm{sh_1}} +(T^{\mathrm{drv}}_{\mathrm{sh_2}} - T^{\mathrm{drv}}_{\mathrm{sh_1}}) z_i ,
    \qquad \, i \in \mathcal{I}.
    \label{eq:sd_ub}
\end{equation}
\begin{alignat}{2}
    z_{i+1} &\geq z_i - \rho_i, &\qquad& i \in \mathcal{I}\setminus\{N\}
    \label{eq:z_prop}\\
    q_i &\geq z_i + \rho_i - 1, &\qquad& i \in \mathcal{I}
    \label{eq:q_def}\\
    \sum_{i \in \mathcal I} q_i + z_N &\leq \bar\epsilon &&
    \label{eq:ext_budget}
\end{alignat}




\subsubsection{Working regulation}
Working time comprises driving, customer service, CS queuing, maneuvering, and charging time not covered by a synchronized break. The work contribution of charging is simply $\tau^c_i - g_i$, which is linear, so no auxiliary variable is needed. The accumulated shift working time $s_i^w$ is reset only by a daily rest.By defining $l^4_i = \rho_i s^w_i$,  constraints~\eqref{eq:l4_ub1}--\eqref{eq:l4_lb} provide the linearization with big-M $T^{spr}_2 = 15$\,h (within a shift, working time is bounded only by the spread). Constraint~\eqref{eq:sw_prop} propagates $s^w_i$. The counter carries no daily cap: the daily-level bound on on-duty time is the shift spread \eqref{eq:h_prop}--\eqref{eq:h_ub} below.

\begin{equation}
    s_{i+1}^w = s_i^w - l_i^4 + D_i
    + \begin{cases}
         M_{i+1}^{stop} v_{i+1} + Q_{i+1}\, y_{i+1} + (\tau^c_{i+1} - g_{i+1}) + M_{i+1}^{seq} \sigma_{i+1}  & \text{if } i+1 \in \mathcal{S}, \\
        S_{i+1}                      & \text{if } i+1 \in \mathcal{C},
    \end{cases}
    \quad \, i \in \mathcal{I}\setminus\{N\}
    \label{eq:sw_prop}
\end{equation}


The shift spread $h_i$ measures the elapsed time since the end of the last daily rest. Since a daily rest is always the last activity performed at a stop, the elapsed time from the end of the previous rest to the start of a rest at stop $i$ is exactly $h_i + o_i$, where $o_i = t^d_i - t^a_i - \tau^r_i$ is the on-duty stop time. Constraint~\eqref{eq:h_prop} propagates the spread, with the reset $l^5_i = \rho_i (h_i + o_i)$ linearized in \eqref{eq:l5_1}--\eqref{eq:l5_3}. Constraint~\eqref{eq:spread_cap} bounds the pre-rest spread at $T^{spr}_1 = 13$\,h when the upcoming rest is regular and $T^{spr}_2 = 15$\,h when it is reduced, implementing the requirement that the daily rest be completed within 24\,h of the end of the previous one. Constraint~\eqref{eq:h_ub} is the global bound for shift length.
\begin{equation}
    h_{i+1} = h_i + o_i + D_i - l^5_i, \qquad i \in \mathcal I \setminus \{N\}
    \label{eq:h_prop}
\end{equation}

\begin{alignat}{2}
    h_i + o_i &\leq T^{spr}_1 + (T^{spr}_2 - T^{spr}_1)\,\rho^{r_2}_i + T^{spr}_2\,(1 - \rho_i), &\qquad& i \in \mathcal I \label{eq:spread_cap}\\
    h_i &\leq T^{spr}_2, &\qquad& i \in \mathcal I \label{eq:h_ub}\\
\end{alignat}



\subsubsection{Initial Conditions}
Constraints~\eqref{eq:init_t_e}--\eqref{eq:init_sw_phi} fix the state at the origin: the truck departs at a given time $t^0$, with initial SOC $E^0$, and all HoS accumulators and the split-break flag are initialized to zero. No break or rest activity is permitted at the origin or at the destination: $x_0 = \rho_0 = x_N = \rho_N = 0$.


\begin{alignat}{2}
    t_0^a &= t^0, &\quad e_0^a &= E^0, \label{eq:init_t_e}\\
    c_0^d &= 0,   &\quad s_0^d &= 0,   \label{eq:init_cd_sd}\\
    s_0^w &= 0,   &\quad \phi_0 &= 0.  \label{eq:init_sw_phi}
\end{alignat}



\begin{figure}
    \centering
    \includegraphics[width=1\linewidth]{figures/ExampleSolution3.png}
    \caption{A feasible schedule on a 950 km corridor with three customers. The four panels share a common time axis. }
    \label{fig:solution}
\end{figure}



{\color{red}
\section{Complexity of the problem (to be removed)}
\subsection{existing work}




SPPRC on general graphs: NP-hard \cite{dror_note_1994} \\
SPPRC on a DAG: Polynomial \cite{irnich_shortest_2005}

\subsection{Charging subproblem}
\noindent No breaks/rests are considered here. Let's consider:
\begin{itemize}
    \item a vector $y \in \{0,1\}^{|\mathcal K|}$ that defines the subset of CS stops $K(y)=\{i \in \mathcal K : y_i =1, \tau_i^c >0\}$ where charging is forced
    \item Non-decreasing convex PWL charging function such that $\tau_i^c = T(e^d_i)-T(e^a_i)$ where T is the PWL function $T:[0,E_{cap}]-> \mathbb{R}$
    \item no condition for final SOC exists
    \item we aim to minimize route duration
\end{itemize}

We aim to prove that: The optimal strategy is to charge each time just enough to reach the next CS in $K(y)$.

The subproblem with fixed charging stops and no HoS constraints is now:
\begin{equation}
    \min t^a_N=t^d_0 + \sum_{i=0}^{N-1}D_i+\sum_{k \in A(y)}(q_k+T(e^d_k)-T(e^a_k)) \equiv \min \sum_{k \in A(y)}(T(e^d_k)-T(e^a_k))
\end{equation}
s.t energy feasibility constraints.


We index the stops in $A(y)$ in route order: $A(y)=\{k_1 < k_2 < ... <k_m\}$ and set $k_{m+1}=N$
we can denote $ j \in [1,m]:\quad C_j = \sum_{l=k_j}^{k_{j+1}-1}E_l > 0$ (energy needed to travel from from this CS to next activated CS, considering all intermediate non activated stops).
We then get $ j \in [1,m],\quad e^a_{k_{j+1}} = e^d_{k_j}-C_j$.
We also define the energy floor at CS departure as the minimum SOC that is required to reach next CS where we can charge: $E_{k_j}^{fl}= E_{min} + C_j, \quad j \in [1,m]$

We can now reformulate the obj as:
\begin{equation}
     \min -T(e^a_{k_1})+\sum_{j=1}^{m-1}(T(e^d_{k_j})-T(e^d_{k_j}-C_{j}))+T(e^d_{k_m})
\end{equation}
where $-T(e^a_{k_1})$ is a constant.

we want to prove that the optimal solution to the problem is $e^d_{k_j}=E^{fl}_{k_j}\quad j \in [1,m]$


{\color{red}TBD, something about using the derivative of the obj function using convexity and non-decreasing argument of the PWL charging function T ?}

\subsection{HoS subproblem}
\noindent No charging stops are considered here. Let's consider:
\begin{itemize}
    \item a vector $x \in \{0,1\}^{|\mathcal K|}$ that defines the subset of CS stops $K(x)=\{i \in \mathcal I : x_i+\rho_i=1, \tau_i^b >0 \text{ or }\tau_i^r > 0\}$ where breaking or resting is forced
    \item no condition for final HoS state exists
    \item we aim to minimize route duration
\end{itemize}

To show: The optimal strategy is to make the shortest possible break or rest each time, while keeping future feasibility of route. Ex: no 15min break possible if we drove more than 4h already, because then not possible to fit a 30min break in the total 4.5h timeperiod for driving.

\noindent \citep{garaix_label-setting_2024}
\begin{itemize}
    \item BUT they considered night work with rolling 24h window that we dont include
    \item they considered break/rest possible anywhere and in the middle of actions
    \item -> they dont prove but suggest that TDSP is NP-complete
\end{itemize}


\noindent \citep{drexl_labelling_2010}
\begin{itemize}
    \item heuristic labelling algorithms for considering EU drivers rules in shortest path problems with resource constraints
    \item Input: An instance of the elementary shortest path problem with time windows and a feasible solution p. Question: Is there a legal schedule for p? Conjecture: Decision problem LEGAL-ROUTE-DR is NP-complete.
\end{itemize}

\noindent \citep{xu_solving_2003}
\begin{itemize}
    \item First paper to model US HoS in a scheduling context (rich pickup-and-delivery with multiple time windows).
    \item  Conjecture: determining a minimum-cost schedule for a given route under US HoS with multiple time windows is NP-hard.
\end{itemize}




\subsection{Combining subproblems}

need to consider $u_i=\tau_i^c(1-x_i-\rho_i)$ which measure the charging time not counted as break/rest.

\noindent options:
\begin{enumerate}
    \item Stops where we charge is fixed, stops where we break/rest is fixed and we combine individual subproblems
    \item We know the stops but without knowing what we do there (break/rest OR charge OR both)
\end{enumerate}

For 1. synchronize if possible, otherwise use decomposition of the two subproblems. Who is dominant to get to next stop when we are at a stop where we both charge and break?

For 2. we have a problem with heteregenous queue duration at CS (not equivalent, forced fixed cost if charging activated at a stop). NP-hard problem ?

Focusing on a singular stop: already hard due to "seuil" problems (cant charge more than 100\%, etc.) + combinatoire entre arrets

}




\section{Solution methods} \label{sec:method}
Introducing uncertainty into a deterministic scheduling model is non-trivial: the optimal schedule depends on travel times and energy consumption that are not yet known at departure. We compare various method to tackle the problem at hand. First a simulation framework is created in Section \ref{sec:simulation} to track the real life progress of a truck, based on some decisions at each stop and realized uncertainty at each leg.
The remainder of this section is organized by \textit{when} a method commits to a decision and \textit{what} information it uses to do so, as summarized in Table \ref{tab:methodologies}: Section \ref{sec:RHLA} and Section \ref{sec:greedy} describe the two adaptive policies, which re-decide at every stop as the trip unfolds, Section \ref{sec:ro}, {Section \ref{sec:robu},} and Section \ref{sec:sp} describe the {three} anticipative methods, which fix a full schedule before departure. Finally, Section \ref{sec:oracle} introduces the perfect hindsight model that take the realized uncertainty as nominal value and solve the equivalent deterministic MILP as a benchmark for the other methods.


\begin{table} [h!]
    \centering
    \begin{tabular}{lll}\toprule
         Methodology&  Decision setting& Travel time distribution\\\midrule
         Rolling-Horizon (LA) &  Online & Scenarios\\
         Greedy&  Online & Average\\
         Robust opt. (RO)&  Offline & Robust over an interval set\\
         Budgeted robust opt. (ROBU)& Offline & Robust over a budgeted set\\
         2 stage SP&  Offline with online recourse & Scenarios\\
         Oracle&  Posteriori& Realized values\\ \bottomrule
    \end{tabular}
    \caption{Solution method characteristics}
    \label{tab:methodologies}
\end{table}

\subsection{Simulation Framework} \label{sec:simulation}
We decouple the planning problem from the execution problem through a simulation framework, illustrated in Figure~\ref{fig:simulation_decision_selector}.

The simulation proceeds stop by stop from origin~$0$ to destination~$N$. At each stop~$i$, the vehicle has arrived with a known state $\Gamma_i = (t^a_i,\, e^a_i,\, c^d_i,\, s^d_i,\, s^w_i,\, h_i,\, \phi_i,\, n^{r2}_i, n^{\epsilon}_i)$ that summarizes accumulated time, energy, and HoS counters, where $n^{r2}_i$ and $n^{\epsilon}_i$ count the reduced rests and driving extensions already consumed. A decision method maps this state to an action $a_i$ specifying whether to charge, which break type to take, and which rest type to take. The vehicle executes $a_i$, departs, and the travel time on leg $i \to i{+}1$ is revealed as a random draw, from which the realized energy consumption follows via the ECR curve. The truck state is updated accordingly and the cycle repeats. Early arrival at a customer triggers immediate service, with the out-of-window penalty recorded

\begin{figure} [h!]
    \centering
    \includegraphics[width=1\linewidth]{figures/simulation.png}
    \caption{The simulation framework. At each stop the realized state $\Gamma_i$ is passed to a decision selector, which returns an action $a_i$ (charge / break type / rest type). The action is executed, the vehicle departs, and only then is the travel-time multiplier $\xi_i$ of the next leg drawn; the realized energy follows from the ECR curve at the implied speed. All six methods share this loop and this state representation, they differ only in the box labeled ``decision selector''.}
    \label{fig:simulation_decision_selector}
\end{figure}


The decisions methods are described below and they share the same simulation loop and vehicle state representation, only the decision step differs. The travel-time realizations are recorded throughout and fed to the oracle at the end of each run.

Because uncertainty is revealed only after departure, a leg may retroactively cause a violation that no stop-level action can repair: the consecutive-driving counter may exceed 4.5 h mid-leg or the SOC may fall below acceptable level. The simulator records these events and runs with any violation are marked infeasible.



\subsection{Rolling-Horizon with look-Ahead policy} \label{sec:RHLA}
We adopt a \emph{rolling-horizon} policy in which a new decision is made at each stop, using only the information available at that point and considering a limited horizon to model the future. Let $i \in \{0, \ldots, N-1\}$ denote the current stop and let $\Gamma_i = (t^a_i,\, e^a_i,\, c^d_i,\, s^d_i,\, s^w_i,\, h_i,\, \phi_i,\, n^{r2}_i, n^{\epsilon}_i)$ be the vehicle state at arrival. At each stop the policy executes the following steps, which are summarized in Figure \ref{fig:RH}.

All structurally feasible actions $a \in \mathcal{A}_i$ are enumerated. An action specifies the charge binary $y \in \{0,1\}$ (relevant at CS stops only), a break type $b \in \{\varnothing, b_{45}, b_{15}, b_{30}\}$, and a rest type $r \in \{\varnothing, r_1, r_2\}$. Actions that violate structural feasibility (such as a split-break second part $b_{30}$ without a prior $b_{15}$, or a reduced rest $r_2$ when the budget is exhausted) are excluded. %Actions whose execution could render the next leg infeasible under the worst-case realization of its travel time (not charging when the worst-case energy demand of the remaining distance to the next reachable charger exceeds the current usable SOC, or not resting when the worst-case leg duration would breach the HoS) are additionally pruned to avoid infeasibility and to reduce computation time at stop level.
No probabilistic feasibility guard is applied: an action that would strand the vehicle is not filtered a priori but is scored as infeasible in the scenarios that expose it, and is penalized through $f(a)$.

Then, an horizon end-stop $h > i$ is determined by advancing along the route for a nominal duration of $T_{\mathrm{hor}}$ hours. This end-stop is defined via a simplistic algorithm that defines, based on rest rule, which stop can be attained in the given horizon.

For each travel leg between our current stop $i$ and the horizon end-stop $h$, a set of $S$ travel-time scenarios $\{ \boldsymbol{\xi}^{(k)}\}_{k=1}^{S}$ is drawn independently for legs $[s, h)$. Each multiplier $\xi^{(k)}_i$ is sampled from the appropriate travel time random distribution. The corresponding scenario energy $E^{(k)}_i$ is derived from the ECR curve evaluated at the implied speed.

Afterward, for each candidate action $a \in \mathcal{A}_i$ and each scenario $k$, the sub-problem MILP over the horizon $[s, h]$ is solved with $a$ fixed at the current stop and with travel times $D^{(k)}_i = D_i \xi^{(k)}_i$. The subproblem minimizes arrival time at $h$, subject to the same SOC, HoS, break, rest, and charging constraints as the full-route model. The score of action $a$ is the aggregation of the $S$ sub-problem objectives:
\[
  f(a) \;=\; \frac{1}{S}\sum_{k=1}^{S} t^{a,(k)}_h(a),
\]
with infeasible scenarios penalized by a large constant. {\color{red} Alternative aggregations (worst-case, best-case, CVAR) to be considered }

Finally, the action $a^* = argmin_ {a \in \mathcal{A}_s} f(a)$ is selected, with a tie-breaking rule that favors charging ($y=1$). A final MILP re-solve on $[s, h]$ with $a^*$ fixed and nominal travel times produces the exact activity durations $(\tau^c_i, \tau^b_i, \tau^r_i)$ executed by the vehicle. The actual travel time to next stop is then drawn from the simulation framework and the vehicle state $\sigma_{i+1}$ is updated accordingly.


\begin{figure}[h!]
    \centering
    \includegraphics[width=1\linewidth]{figures/RH.png}
    \caption{Rolling horizon with look ahead policy.}
    \label{fig:RH}
\end{figure}

In order to speed-up the process, only a partially relaxed LP version of the subproblems are solved. The out-of-window indicators $\delta_i$ and the next stops decisions remains as integer or binary, as well as all binary decisions at customer stops and the sequential charging binary at all CS. However, to tighten the relaxation, various valid inequalities are added to the formulation of the LP subproblems. They are summarized in \ref{sec:valid_inequ}.

\subsection{Greedy heuristic} \label{sec:greedy}

To provide a baseline representative of current practice, we implement a \emph{greedy} policy that mimics the behavior of an driver operating without planning tools: stop only when forced to, charge to full whenever stopping, and avoid queuing stations known to be busy. Unlike the rolling-horizon policy, the greedy algorithm uses no look-ahead and solves no optimization sub-problem: it evaluates a fixed priority rule at each stop and returns an action immediately. A break can be inserted in the charging time if the duration is long enough.


\subsection{Robust optimization} \label{sec:ro}
Travel times are uncertain, which we model as
\[
  \tilde{D}_i = D_i \,\xi_i, \qquad \xi_i \in {[\underline\xi,\ \overline\xi]}, \qquad i \in \mathcal I ,
\]
with physically motivated, asymmetric bounds rather than a symmetric perturbation: $\overline\xi = V^{\mathrm{nom}}/\bar v$ from a leg-average congestion floor $\bar v$, and $\underline\xi = V^{\mathrm{nom}}/\hat v$ from the vehicle speed limiter $\hat v$ (90~km/h for HGVs).

Since the ECR curve makes energy speed-dependent, each $\tilde D_i$ also induces an uncertain energy, and the robust counterpart is obtained by parameter substitution at opposite extremes of the interval:
\[
  D_i^{\mathrm{wc}} = D_i\,{\overline\xi}, \qquad
  E_i^{\max} = L_i \cdot \mathrm{ECR}\!\left(\frac{L_i}{D_i\,{\underline\xi}}\right),
  \qquad i = 0,\ldots,N-1.
\]
The slowest realization is unfavorable for elapsed time and the Hours-of-Service accumulators, the fastest for state of charge, since the ECR curve is increasing and convex in speed. All other constraints are unchanged, and the resulting plan is feasible for every realization $\xi_i \in {[\underline\xi, \overline\xi]}$.

This box is the most conservative member of the budgeted family of \citet{bertsimas_price_2004}, recovered at full budget $\Gamma = N$, and we retain it deliberately as the maximum-commitment plan: every activity and every duration is fixed before departure from the support alone, and executed as is in the simulation framework. Section~\ref{sec:robu} relaxes the box to a budgeted set that removes most of this conservatism while retaining a deterministic feasibility guarantee.


\subsection{Budgeted robust optimization} \label{sec:robu}
The box counterpart of Section~\ref{sec:ro} prices every leg at its worst case simultaneously, a joint event of very low probability under independent leg multipliers. We therefore also field the budgeted uncertainty set of \citet{bertsimas_price_2004},
\[
  \mathcal U_\Gamma = \Bigl\{ \xi \;:\; \xi_i = 1 + z_i^{+}(\overline\xi - 1) - z_i^{-}(1-\underline\xi),
  \;\; z^{+}\!, z^{-} \in [0,1]^N,
  \;\; \textstyle\sum_i z_i^{+} \le \Gamma,
  \;\; \sum_i z_i^{-} \le \Gamma \Bigr\},
\]
in which at most $\Gamma$ legs' worth of deviation mass is slow ($\xi_i \to \overline\xi$, unfavorable for time and Hours of Service) and, independently, at most $\Gamma$ legs' worth is fast ($\xi_i \to \underline\xi$, unfavorable for state of charge).
%We use the classical protection level $\Gamma = \lceil 1 + z_{1-\varepsilon}\sqrt{N} \rceil$ with $\varepsilon = 0.01$, for which any single constraint over $N$ independent, symmetrically distributed multipliers is violated with probability at most $\varepsilon$ \citep{bertsimas_price_2004}. The budget interpolates between the nominal model ($\Gamma = 0$) and the box of Section~\ref{sec:ro} ($\Gamma = N$).
We use the classical protection level $\Gamma = \sqrt{N} $ \citep{bertsimas_price_2004}. The budget interpolates between the nominal model ($\Gamma = 0$) and the box of Section~\ref{sec:ro} ($\Gamma = N$).

Unlike the box, the budgeted counterpart does not reduce to parameter substitution: the accumulators (driving since the last break, shift driving, state of charge) reset at stops selected by the binary plan, so each robust constraint involves the worst deviation over a decision-dependent window of legs. Constraint-wise dualization then yields products of dual variables and plan binaries and is impractical at our instance sizes \citep{poss_robust_2013}, an obstacle the vehicle routing literature also meets inside branch-and-price labeling \citep{munari_robust_2019}.
We instead solve the counterpart exactly by scenario generation \citep{mutapcic_cutting-set_2009, zeng_solving_2013}. A master problem optimizes the plan against a finite subset $\widehat{\mathcal U} \subset \mathcal U_\Gamma$ --- the extensive-form model of Section~\ref{sec:sp} with a min--max objective and a single non-anticipative duration vector --- and a pessimization step searches $\mathcal U_\Gamma$ for a realization that breaks the current plan. With the plan fixed every accumulator window is known, so the worst attack on each constraint family is available in closed form: fully deviate the $\Gamma$ largest legs inside that window. Candidates are generated per continuous-driving window, shift, customer-window prefix, and between-charges segment, each certified against the full model; violated scenarios are appended to $\widehat{\mathcal U}$ and the master re-solved. The procedure terminates finitely \citep{zeng_solving_2013}, in practice within a handful of iterations, and the master may be solved to a relative gap without losing convergence \citep{patzold_approximate_2020}.

At termination the plan is provably feasible for every realization in $\mathcal U_\Gamma$ and, as for the box, commits every activity and every duration before departure and is executed as is in the simulation framework. By the calibration of $\Gamma$, realizations outside the set may occur and render it infeasible.


\subsection{Two-stage stochastic program}
\label{sec:sp}

We extend the deterministic MILP into a two-stage stochastic program (2-SP) with fixed recourse. The first stage holds the binary decisions at each stop (charge, break, rest yes/no), fixed before any uncertainty is revealed and hence identical across scenarios;
the second stage holds the activity durations at each stop, free to adapt to each scenario's realized travel times.

Let $k = 1, \ldots, S$ index the scenarios, each a random realization of travel time on every leg, drawn independently across legs (the same sampling assumption as the rolling-horizon sub-problem of Section~\ref{sec:RHLA}). Denoting the first-stage binaries $y, b, r$ and the second-stage variables $\tau^{c,(k)}, \tau^{b,(k)}, \tau^{r,(k)}$, the model minimizes the expected arrival time,
\[
  \min_{y,b,r \,\in\, \{0,1\}} \;\; \frac{1}{S}\sum_{k=1}^{S} \bigl(t^{a,(k)}_N(y,b,r) + \beta \sum_{i \in \mathcal C}\delta_i^{(k)}\bigr)
\]
subject to the SoC, HoS, break, rest, and charging constraints of the deterministic model holding for every scenario $k$, with nominal travel times and energies replaced by their scenario-specific values.

Only the first-stage binaries are retained from the offline solve. During simulation, at each stop $i$ they are held fixed and the durations over $[i, N]$ are re-optimized as a linear program from the realized state $\sigma_i$ and nominal remaining travel times; only the durations at stop $i$ are executed. If the LP is infeasible, a repair step solves a restricted MILP in which binary activities may be added but not removed (for example an extra CS stop), minimizing first the number of additions and then arrival time. If the repair also fails, the run is recorded as a violation. The plan's structure is thus committed offline, its durations adapt online, and its repair frequency is itself a reported metric.



\subsection{Perfect-Hindsight Oracle} \label{sec:oracle}
Unlike the methods above, the oracle is not a candidate policy: it is an ex-post performance bound computed once the simulation run is finished, using information no policy could have had while driving.
To assess how much is lost by the lack of advance knowledge, we compare the methodologies against a perfect-hindsight oracle: a lower bound on achievable arrival time given the travel times that actually occurred during the simulation run.

After the simulation has completed, let $\tilde{D}_0, \tilde{D}_1, \ldots, \tilde{D}_{N-1}$ be the sequence of travel times that were realized on each leg. The oracle solves the full-route MILP of Section~\ref{sec:model} with these values substituted for the nominal times $D_i$. Because the realized travel times are fixed and known when the problem is solved, this is a deterministic MILP identical in structure to the nominal model. Its optimal value $t^{a,*}_N$ is the earliest possible arrival at the destination that any policy could have achieved, had it known all travel times in advance.

The oracle is not a realizable policy: no algorithm can know $\tilde{D}_i$ before driving leg~$i$. It therefore provides a lower bound against which the optimality gap of any online policy can be measured.  The hindsight gap $\Delta$ is the relative difference of the full route duration. A small $\Delta$ indicates that look-ahead scheduling recovers most of the benefit of perfect foresight. In order to speed up the MILP solve, the rolling horizon solution is used as a first feasible solution in the perfect hindsight model.


\section{Experimental design} \label{sec:exp}
This section presents the setting of our experimental setup, in order to provide reproducibility of our results. The full data collection is made available in Github XXX.

\subsection{Nominal parameters}
All experiments share the vehicle and operational parameters listed in Table~\ref{tab:params}.

\begin{table}[ht]
\centering
\caption{Nominal vehicle and operational parameters.}
\label{tab:params}
\begin{tabular}{llr}
\toprule
\textbf{Category} & \textbf{Parameter} & \textbf{Value} \\
\midrule
\multirow{3}{*}{Vehicle}
  & Battery capacity $E^{\mathrm{cap}}$        & 500\,kWh \\
  & Minimum SoC $E^{\min}$                      & $0.2\cdot E^{\mathrm{cap}}$ (100\,kWh) \\
\midrule
\multirow{2}{*}{Charging}
  & Charger power output                      & 350 kW \\
  & PWL breakpoints (SoC)                       & 0\%, 80\%, 100\% \\
\midrule
\multirow{4}{*}{Operations}
  & Nominal average speed                       & 80\,km/h \\
  & Customer service time                       & 30\,min \\
  & Maneuver time when seq. operations at CS stop               & 5\,min \\
  & Maneuver time at CS and rest areas stop               & 10\,min \\
  & Out of window penalty $\beta$ & 30\,min \\
\midrule
\multirow{2}{*}{Infrastructure}
  & Charging station spacing                    & $U[40, 80]$\,km \\
  & Charging station or rest areas spacing      & $<$25\,km \\
  & Queue delay at CS (log-normal)              & $\mu = 10$\,min, $\sigma = 8$\,min \\
\bottomrule
\end{tabular}
\end{table}

Charging station spacing and charger power output is taken from the Regulation (EU) 2023/1804, which mandates heavy-duty vehicles charging station every 60 km on the TEN-T core network by 2030, and advise for charger power higher than 350 kW. The energy consumption rate (ECR) as a function of speed $v$ (km/h) is modeled from \citet{gao_estimating_2025} and calibrated with verified consumption values from constructors.

%Energy consumption follows the simplified road-load formulation of Earl et al. (2018), calibrated for a 40 t tractor-semitrailer: C_RR = 0.0055, frontal area 10 m2, rho = 1.2 kg/m3, drivetrain efficiency 0.85, and a continuous 3.2 kW auxiliary load. We take C_D = 0.50, between the 2018 EU fleet average (0.60) and the best-in-class value (0.36), reflecting the aerodynamic redesign of current long-haul BETs. This yields 1.23 kWh/km at 80 km/h, within the 1.1-1.4 kWh/km range reported for 40 t battery electric tractors

Travel time on each leg is uncertain through a multiplicative multiplier $\tilde D_i = D_i\,\xi_i$. The multiplier follows a shifted-lognormal law $\xi_i = \min\!\bigl(\underline\xi + Y_i,\ \overline\xi\bigr)$, where $Y_i$ is lognormal with a base coefficient of variation of $0.15$, and the fixed bounds are $\underline\xi = V^{\mathrm{nom}}/90 \approx 0.889$ (speed limits) and $\overline\xi = V^{\mathrm{nom}}/50 = 1.6$ (a leg-average congestion floor). These same bounds $[\underline\xi,\overline\xi]$ define the box and budgeted uncertainty sets used by the robust methods of Sections~\ref{sec:ro}--\ref{sec:robu}. The realized energy on each leg follows from the ECR curve at the implied speed.


\subsection{Benchmark instances}

A set of benchmark instances of increasing complexity is generated along a single long-haul corridor. Charging station stops are placed at regular average intervals of 60\,km, customer stops are randomly place around clusters centers inserted along the corridor. The instances are parameterized based on route distance, customer count, and time windows width as summarized in Table~\ref{tab:instances}. For each instance, service windows are generated from the deterministic MILP solved with average travel time: each window is centered on the service start of the found optimal solution, with half-width drawn independently for each customer from the class range in Table \ref{tab:instances}. Instances whose deterministic model is infeasible under nominal data are discarded and regenerated. Each Route/Customers/Time Windows class contains 25 independent and randomly generated instances. The number of stops averages 50 on the short routes, 96 on the medium and 167 on the long ones, of which 16, 32 and 58 are charging stations and 26, 56 and 101 are rest-area nodes.

\begin{table}[ht]
\centering
\caption{Experimental design. Full factorial,
         $3 \times 3 \times 4 \times 25 = 900$ instances.}
\label{tab:instances}
\small
\begin{tabular}{lll}
\toprule
\textbf{Factor} & \textbf{Levels} & \textbf{Values} \\
\midrule
Route class   & Short  & 800--1\,200\,km     \\
              & Medium & 1\,500--2\,500\,km  \\
              & Long   & 3\,000--4\,000\,km  \\
\midrule
Customers     & Few    & 1--3    \\
              & Medium & 4--6    \\
              & Many   & 7--15   \\
\midrule
Time window   & None   & no window         \\
 (half-width  & Tight  & $U[0.5, 1.0]$\,h  \\
 per customer)& Medium & $U[1.0, 3.0]$\,h  \\
              & Large  & $U[3.0, 6.0]$\,h  \\
\bottomrule
\end{tabular}
\end{table}

Each route instance class targets a distinct operational regime: Small instances complete with a single rest, Medium instances span two to three shifts, activating the rest, reduced-rest, and extension mechanisms. Large instances approach the weekly driving cap.

\subsection{Real-life case study}
\label{sec:casestudy}

We complement the synthetic corridors with one instance reconstructed from operational data: a six-day international tour driven by a Norwegian haulier in September~2025, recorded by the vehicle at one-minute resolution (carrier and customers anonymised).
The stop sequence and odometer allow to reconstruct the $3\,339$\,km route, from southern Norway to the central Netherlands and north-eastern Bavaria and back, with a sea crossing in each direction, with nominal leg times taken from the speeds observed on each leg.
Seven customers are identified from the freight order register, with their observed service times ($0.5$--$2.0$\,h, mean $1.2$\,h).
Each ferry crossing collapses into a single rest area node with $x^{b45}_i=1$ and $\tau^b_i$ fixed to the measured port-to-port duration ($3.8$\,h and $4.9$\,h), which requires no new constraints: it never enters the driving or working accumulators, and because the reset indicator already contains $x^{b45}$ it resets the continuous-driving clock.
%the driver's breaks and rests are discarded, since when to stop is what the model decides.
Charging and rest infrastructure come from OpenStreetMap: $|\mathcal K|=43$ nodes restricted to heavy-vehicle accessible sites (150--1\,400\,kW, median 400\,kW), and $|\mathcal L|=63$ rest areas, with off-route sites charged a maneuver time for the detour in excess of a normal stop.
Battery, charging curve, ECR and hours-of-service parameters are unchanged from Table~\ref{tab:params}, and the largest gap between consecutive chargers, $293$\,km, leaves the instance feasible on the base $500$\,kWh pack.
At $3\,339$\,km, seven customers and five duty cycles the instance belongs to the Long/Many class of Table~\ref{tab:instances}, though the real corridor is sparser than the $60$\,km charger spacing the synthetic instances assume. Figure \ref{fig:realroute} illustrate the route instance characteristics.
We assign no time windows, so the case study reads against the Long/Many/None row of Table~\ref{tab:gap}.
\begin{figure}
    \centering
    \includegraphics[width=1\linewidth]{figures/route_real.png}
    \caption{{\color{red}put in appendix ? upgrade and improve anyway}}
    \label{fig:realroute}
\end{figure}


\section{Computational results} \label{sec:results}
The experiments reported in this section have been solved using Gurobi 12.0.2 on a computer equipped with an Intel(R) Xeon(R) Gold 6244 CPU @ 3.60GHz processor and 384 GB of RAM.


\subsection{Validation of the deterministic model on TDSP benchmark instances}
\label{sec:validation}

Before turning to the stochastic experiments, we verify that the HoS core of the deterministic MILP of Section~\ref{sec:model} reproduces known optimal solutions of the pure truck driver scheduling problem. To this end we use the 40 R15-PGLT benchmark instances made available by \citet{garaix_label-setting_2024}\footnote{\url{https://emse.fr/~garaix/TruckDriverSchedule/Expert_Systems/}}, which comprise fixed customer sequences with 2--14 customers, per-stop service times (including a loading operation at the origin depot), hard time windows, and all durations rounded to multiples of 15~minutes. We compare against the published optima of their \emph{no-night} configuration, which enforces exactly the daily provisions of Regulation~(EC) No~561/2006 and Directive 2002/15/EC considered in this paper, without the night-work rule.

Since these are diesel instances, the benchmark exercises the HoS machinery of our model in isolation: we disable the battery ($E_i = 0$,$\mathcal K = \emptyset$) and reduce all maneuver, queue, and out-of-window penalty terms to zero, with time windows enforced as hard constraints ($\delta_i = 0$). Furthermore, the reference model allows breaks and rests to preempt driving and service at arbitrary points. We emulate this on the 15-minute by inserting rest-area nodes ($\mathcal L$) every 15~minutes of driving and by splitting each service operation into 15-minute segment.

%Five differences between the two formulations are bridged as follows.
%First, the reference model allows breaks and rests to preempt driving and service at arbitrary points. We emulate this on the 15-minute by inserting rest-area nodes ($\mathcal L$) every 15~minutes of driving and by splitting each service operation into 15-minute segments, the time window applying to the start of the first segment and the whole operation having to complete by the closing of the window plus its duration.
%Second, a zero-length rest-area node inserted before each customer permits the arrive--rest--serve pattern present in several reference optima.
%Third, their schedules must terminate with a daily rest at the destination (regular, or reduced against the weekly budget $\bar\rho$), whose duration counts toward the objective. We add the corresponding constraint $\rho^{r_1}_{j} + \rho^{r_2}_{j} = 1$ at the final stop.
%Fourth, the departure time from the origin is left free (bounded below by the earliest start), as some instances open the depot window after time zero.


Table~\ref{tab:validation} summarizes the outcome. Our formulation, solved with Gurobi to zero optimality gap, reproduces the published optimal objective exactly on 39 of the 40 instances, with solve times below 64~s (median below 1~s) against instance sizes of up to 180 stops after discretization. On the single divergent instance (Test~4) our model attains 1410~min against a published value of 1425~min.

\begin{table}[htbp]
\centering
\caption{Validation on the 40 R15-PGLT instances (no-night configuration).}
\label{tab:validation}
\small
\begin{tabular}{lr}
\toprule
Instances solved to proven optimality & 40/40 \\
Exact matches & 39/40 \\
Divergent instances & 1 (ours 1410 vs.\ 1425, see text) \\
Stops after 15-min discretization & 23--180 \\
Solve time (median / max) & $<1$~s / 64~s \\
\bottomrule
\end{tabular}
\end{table}


\subsection{Base-case comparison}
\label{sec:basecase}

Each instance class is instantiated with 25 independent random realizations, and every realization is solved by all five methods (Greedy, RO, ROBU, LA, SP) and evaluated in the common discrete-event simulator of Section~\ref{sec:simulation} against the perfect-hindsight oracle.
Every offline optimisation model (RO, ROBU, the two-stage stochastic program, and the oracle) was solved with Gurobi under a gap tolerance of 0.5\% and a time limit of 7,200s. LA's look-ahead subproblems are LP relaxations with a 300s cap that is never binding in practice. Greedy uses no solver.

We compare the methods along three metrics. First, efficiency, shown through the gap to oracle in realized route duration in Table~\ref{tab:gap} and Figure~\ref{fig:resultsgap}. Then, reliability, via the rate at which the executed schedule becomes infeasible (the infeasibility strip of Figure~\ref{fig:resultsgap} and Table~\ref{tab:feasibility}). Finally, computational effort in Table~\ref{tab:runtime}.


Because the oracle is itself a mixed-integer program, the reference against which every method gap is measured is only as good as the oracle gap.
On short and medium routes the oracle closes to optimality for most instancces. However, on \blue{about a third of} long routes the incumbent stabilizes early while the dual bound stalls at \blue{a mean of 6.9\,\% and at most 9.2\,\%}. Long-route gaps in Table~\ref{tab:gap} are therefore not exact distances to the proved optimum, but only to the best incumbent.
\green{Comment: ``on half of long routes'' was too pessimistic against the current
caches. Over the 100 long-route oracle solves logged in
data\_output/additional\_diesel\_stats.csv, 35\,\% end above the 0.5\,\% tolerance,
those ones at 6.9\,\% on average and 9.2\,\% at worst; short is 0/100 and medium 1/100.
The honest phrasing is ``certified to 9\,\%'' on the long class only.}

\green{Comment: the Medium and Large time-window rows of Table~\ref{tab:gap}, blank in
the version you sent, are now largely filled: completely on the short routes, and on
the medium and long routes for Greedy and LA everywhere plus RO and SP on the
\emph{Few} rows. What is still missing is Medium/Large under RO, ROBU and SP on the
medium and long \emph{Medium}/\emph{Many} rows, and the long \emph{Medium} customer
rows entirely. Because the two new window classes are much easier than Tight, adding
them pulls several of the aggregate statements below downwards --- see the
window-penalty paragraph --- and the caption range widens to 12--25 realizations,
the 12 coming from the long/Many/Medium Greedy cell.}

% -----------------------------------------------------------------------------
%  GAP TABLE -- from tex/tables/gap.tex (2026-08-14 20:51), ONE CELL OVERRIDDEN.
% -----------------------------------------------------------------------------
\green{Comment on the table below: it is tex/tables/gap.tex verbatim except for the
RO / Short / Many / None cell. The generated file gives 4.2\,\%, but gap\_nopen.tex
shows that cell aggregated a single run ($n=1$), whereas
data\_output/paper\_gap\_stats.csv --- written in the same pass --- gives 27.6\,\% over
$n=25$. Every other RO cell lies between 21 and 41\,\%, so 4.2 is an artefact of that
one-run pass and I have kept your 27.6. This is the same discrepancy as in the
previous revision, so the table generator still has a stale-run problem worth
fixing before submission.}

\begin{table}[htbp]
\centering
\caption{Gap to oracle (\%) in route duration (window penalties included) by instance class and method, averaged over \blue{12--25} realizations per instance class.  Parentheses give the same gap over the same runs with window penalties excluded (route duration only).  ROBU is solved on the short-route instances only.  LA uses 25 scenarios over a 24\,h horizon; SP uses 25 scenarios.}
\label{tab:gap}
\resizebox{\textwidth}{!}{%
\begin{tabular}{lll c ccccc}
\toprule
\multirow{2}{*}{\textbf{\small Route}} & \multirow{2}{*}{\textbf{\small Cust.}} & \multirow{2}{*}{\textbf{\small TW}} & \multirow{2}{*}{\textbf{Oracle (h)}} &
\multicolumn{5}{c}{\textbf{Gap to Oracle \%}} \\
\cmidrule(lr){5-9}
 & & & & Greedy & RO & ROBU & LA & SP \\
\midrule
\multirow{12}{*}{\small Short} & \multirow{4}{*}{\small Few} & {\small None} & 26.7 & \blue{5.5~{\scriptsize(5.5)}} & 22.3~{\scriptsize(22.3)} & 17.1~{\scriptsize(17.1)} & 5.0~{\scriptsize(5.0)} & 2.1~{\scriptsize(2.1)} \\
 &  & {\small Tight} & 26.1 & \blue{7.5~{\scriptsize(5.7)}} & 23.0~{\scriptsize(20.1)} & 13.0~{\scriptsize(10.4)} & 6.6~{\scriptsize(5.7)} & 2.5~{\scriptsize(2.5)} \\
 &  & {\small Medium} & \blue{26.1} & \blue{7.0~{\scriptsize(6.2)}} & \blue{21.6~{\scriptsize(19.2)}} & -- & \blue{6.2~{\scriptsize(5.4)}} & \blue{2.2~{\scriptsize(2.1)}} \\
 &  & {\small Large} & \blue{26.4} & \blue{6.7~{\scriptsize(5.9)}} & \blue{21.8~{\scriptsize(20.3)}} & -- & \blue{5.8~{\scriptsize(5.0)}} & \blue{2.1~{\scriptsize(2.1)}} \\
 & \multirow{4}{*}{\small Medium} & {\small None} & 28.0 & \blue{5.6~{\scriptsize(5.6)}} & 21.7~{\scriptsize(21.7)} & 17.7~{\scriptsize(17.7)} & 5.3~{\scriptsize(5.3)} & 2.2~{\scriptsize(2.2)} \\
 &  & {\small Tight} & 28.2 & \blue{11.2~{\scriptsize(5.9)}} & 31.4~{\scriptsize(25.1)} & 23.7~{\scriptsize(17.4)} & \blue{8.8~{\scriptsize(5.6)}} & 3.5~{\scriptsize(2.5)} \\
 &  & {\small Medium} & \blue{27.2} & \blue{8.1~{\scriptsize(6.2)}} & \blue{24.4~{\scriptsize(19.9)}} & -- & \blue{6.9~{\scriptsize(5.7)}} & \blue{4.3~{\scriptsize(4.1)}} \\
 &  & {\small Large} & \blue{28.2} & \blue{6.9~{\scriptsize(5.5)}} & \blue{33.5~{\scriptsize(29.5)}} & -- & \blue{6.8~{\scriptsize(5.5)}} & \blue{2.7~{\scriptsize(2.5)}} \\
 & \multirow{4}{*}{\small Many} & {\small None} & 31.0 & 5.3~{\scriptsize(5.3)} & \red{27.6}~{\scriptsize(\red{27.6})} & 22.4~{\scriptsize(22.4)} & 3.9~{\scriptsize(3.9)} & 2.1~{\scriptsize(2.1)} \\
 &  & {\small Tight} & 31.4 & \blue{15.3~{\scriptsize(5.8)}} & 41.1~{\scriptsize(28.2)} & 36.9~{\scriptsize(24.3)} & 8.8~{\scriptsize(4.5)} & 5.7~{\scriptsize(4.3)} \\
 &  & {\small Medium} & \blue{31.5} & \blue{6.9~{\scriptsize(4.6)}} & \blue{36.8~{\scriptsize(28.4)}} & -- & \blue{7.0~{\scriptsize(5.1)}} & \blue{3.6~{\scriptsize(3.3)}} \\
 &  & {\small Large} & \blue{30.0} & \blue{7.9~{\scriptsize(5.2)}} & \blue{32.4~{\scriptsize(26.1)}} & -- & \blue{7.9~{\scriptsize(5.8)}} & \blue{1.8~{\scriptsize(1.7)}} \\
\midrule
\multirow{12}{*}{\small Medium} & \multirow{4}{*}{\small Few} & {\small None} & 54.8 & 6.0~{\scriptsize(6.0)} & 38.1~{\scriptsize(38.1)} & -- & 5.1~{\scriptsize(5.1)} & 6.7~{\scriptsize(6.7)} \\
 &  & {\small Tight} & 56.8 & \blue{6.3~{\scriptsize(5.0)}} & 30.4~{\scriptsize(28.9)} & -- & 5.4~{\scriptsize(4.4)} & 4.3~{\scriptsize(3.8)} \\
 &  & {\small Medium} & \blue{54.3} & \blue{6.1~{\scriptsize(5.6)}} & \blue{32.3~{\scriptsize(31.0)}} & -- & \blue{5.6~{\scriptsize(5.3)}} & \blue{2.5~{\scriptsize(2.5)}} \\
 &  & {\small Large} & \blue{55.4} & \blue{7.4~{\scriptsize(7.1)}} & \blue{36.4~{\scriptsize(34.6)}} & -- & \blue{7.6~{\scriptsize(7.3)}} & \blue{5.2~{\scriptsize(5.1)}} \\
 & \multirow{4}{*}{\small Medium} & {\small None} & 58.0 & 5.0~{\scriptsize(5.0)} & 30.2~{\scriptsize(30.2)} & -- & 4.5~{\scriptsize(4.5)} & 3.3~{\scriptsize(3.3)} \\
 &  & {\small Tight} & 56.4 & \blue{8.5~{\scriptsize(5.5)}} & 33.1~{\scriptsize(29.1)} & -- & 7.1~{\scriptsize(4.5)} & 4.8~{\scriptsize(3.5)} \\
 &  & {\small Medium} & \blue{55.6} & \blue{7.1~{\scriptsize(5.1)}} & -- & -- & \blue{6.3~{\scriptsize(4.9)}} & -- \\
 &  & {\small Large} & \blue{54.7} & \blue{6.7~{\scriptsize(5.7)}} & -- & -- & \blue{6.0~{\scriptsize(5.3)}} & -- \\
 & \multirow{4}{*}{\small Many} & {\small None} & 60.6 & \blue{6.3~{\scriptsize(6.3)}} & 28.6~{\scriptsize(28.6)} & -- & 6.7~{\scriptsize(6.7)} & 4.9~{\scriptsize(4.9)} \\
 &  & {\small Tight} & 59.8 & \blue{13.5~{\scriptsize(5.3)}} & 37.0~{\scriptsize(28.2)} & -- & 14.8~{\scriptsize(7.4)} & 7.7~{\scriptsize(4.0)} \\
 &  & {\small Medium} & \blue{59.4} & \blue{10.1~{\scriptsize(6.5)}} & -- & -- & \blue{7.9~{\scriptsize(5.5)}} & -- \\
 &  & {\small Large} & \blue{60.2} & \blue{8.3~{\scriptsize(6.9)}} & -- & -- & \blue{6.5~{\scriptsize(5.2)}} & -- \\
\midrule
\multirow{12}{*}{\small Long} & \multirow{4}{*}{\small Few} & {\small None} & 104.3 & \blue{5.4~{\scriptsize(5.4)}} & 35.2~{\scriptsize(35.2)} & -- & 5.0~{\scriptsize(5.0)} & 10.0~{\scriptsize(10.0)} \\
 &  & {\small Tight} & 104.2 & \blue{6.2~{\scriptsize(5.4)}} & 35.2~{\scriptsize(34.3)} & -- & 4.9~{\scriptsize(4.5)} & 11.4~{\scriptsize(10.7)} \\
 &  & {\small Medium} & \blue{105.6} & \blue{7.0~{\scriptsize(6.6)}} & -- & -- & \blue{9.3~{\scriptsize(9.0)}} & \blue{10.9~{\scriptsize(10.4)}} \\
 &  & {\small Large} & \blue{104.1} & \blue{5.8~{\scriptsize(5.7)}} & -- & -- & \blue{5.7~{\scriptsize(5.6)}} & \blue{5.9~{\scriptsize(5.7)}} \\
 & \multirow{4}{*}{\small Medium} & {\small None} & 102.0 & \blue{5.3~{\scriptsize(5.3)}} & 32.4~{\scriptsize(32.4)} & -- & 5.1~{\scriptsize(5.1)} & 8.2~{\scriptsize(8.2)} \\
 &  & {\small Tight} & 104.2 & \blue{7.9~{\scriptsize(6.1)}} & \blue{37.7~{\scriptsize(35.5)}} & -- & 7.2~{\scriptsize(5.8)} & 12.0~{\scriptsize(10.3)} \\
 &  & {\small Medium} & \blue{--} & \blue{--} & \blue{--} & -- & \blue{--} & \blue{--} \\
 &  & {\small Large} & \blue{--} & \blue{--} & \blue{--} & -- & \blue{--} & \blue{--} \\
 & \multirow{4}{*}{\small Many} & {\small None} & 108.6 & \blue{7.0~{\scriptsize(7.0)}} & 32.7~{\scriptsize(32.7)} & -- & 5.1~{\scriptsize(5.1)} & 11.8~{\scriptsize(11.8)} \\
 &  & {\small Tight} & 103.6 & \blue{9.5~{\scriptsize(6.3)}} & 36.6~{\scriptsize(31.7)} & -- & 7.7~{\scriptsize(5.2)} & 17.2~{\scriptsize(13.9)} \\
 &  & {\small Medium} & \blue{100.6} & \blue{8.2~{\scriptsize(7.4)}} & -- & -- & \blue{5.3~{\scriptsize(4.5)}} & -- \\
 &  & {\small Large} & \blue{107.3} & \blue{6.6~{\scriptsize(5.7)}} & -- & -- & \blue{6.6~{\scriptsize(6.0)}} & -- \\
\bottomrule
\end{tabular}%
}
\end{table}


\begin{figure}[htbp]
    \centering
    \includegraphics[width=\linewidth]{figures/resultsgappooledTW.png}
    \caption{Distribution of the realized gap to the oracle across seeds, by route class, customer count, method (colour) and aggregated over time-window class.  Boxes span the interquartile range, whiskers 1.5\,IQR.  The heat strip below each panel shades the infeasibility rate. Empty slots are instance/method combinations without completed
    runs.}
    \label{fig:resultsgap}
\end{figure}



Across short and medium instance classes the two-stage stochastic plan (SP) and the look-ahead policy (LA) attain the smallest gaps to the oracle, as seen in Figure \ref{fig:resultsgap} and Table \ref{tab:gap}. The myopic greedy heuristic is intermediate, while both static robust plans are the most conservative.
This ordering follows the amount of information each method uses and the point at which it commits: the two policies that adapt to realized travel times (SP through online duration recourse and LA through full re-optimisation at every stop) recover most of the hindsight optimum, whereas the offline robust plans pay for committing every duration in advance against a possible worst case.
The ordering between SP and LA reverses with route length: on the short routes SP is clearly ahead (\blue{2.9\%} vs.\ \blue{6.6\%} on average) because 25 scenarios over the whole horizon still capture the realization well, while on the medium routes the two are close (\blue{5.1\%} and \blue{6.9\%}) and on the long routes SP terminates due to time limit because of its exponential size, meaning large optimality gap and therefore potential worse solutions given to the simulation framework. Greedy, conversely, improves in relative terms as routes lengthen.

The gap widens sharply with the number of customers and the tightness of the windows. On the short routes the \emph{Few}/\emph{None} cell against the \emph{Many}/\emph{Tight} show a steep deterioration for the three committing methods and less than a doubling for the two adaptive ones. A static schedule that must satisfy more, and narrower, service windows simultaneously is exactly the situation in which committing early is most expensive.
\green{Comment: ``less than a doubling for the two adaptive ones'' no longer separates
the methods cleanly. Short/Few/None to Short/Many/Tight is
$5.0 \to 8.8$ for LA ($1.8\times$) but $2.1 \to 5.7$ for SP ($2.7\times$), while RO only
goes $22.3 \to 41.1$ ($1.8\times$) --- the same factor as LA. The multiplicative reading
does not hold; what does hold is the additive one, and I would rephrase as
``costs the committing methods 14--19 percentage points and the adaptive ones
under 4''.}

A substantial part of the gap on the window-constrained classes is window penalty rather than slower routing. Comparing the penalized gap of Table~\ref{tab:gap} with the duration-only gap in parentheses, the averaged $\beta\sum_i\delta_i$ term accounts for \blue{22\,\%} of Greedy's gap, \blue{18\,\%} of LA's, \blue{12\,\%} of ROBU's, \blue{10\,\%} of SP's and only \blue{9\,\%} of RO's. The ordering is informative: RO is genuinely slower, spending its gap on conservative durations rather than on missed windows, while Greedy's tight-window gap is mostly the cost of arriving off the windows: on short/Many/Tight, \blue{15.3\,\%} penalised against 5.8\,\% duration-only, so almost two thirds of it. Read on duration alone, Greedy is close to LA everywhere.
\green{Comment: the ordering you rely on (Greedy $>$ LA $>$ ROBU $>$ SP $>$ RO) survives,
but every share dropped by 10--15 points against your version, for two reasons worth
knowing. First, the Medium and Large window classes are now in the average and they
carry almost no penalty, which dilutes it: over None+Tight alone the same shares are
25/20/12/11/8\,\%. Second, the greedy fix changed the mix --- Greedy now misses
fewer windows relative to how much duration it loses. Note also that SP and RO have
converged to within a point of each other (10 against 9), so the tail of the
ordering no longer carries any signal; I would name only Greedy and LA at one end
and RO at the other.}

Where both static robust variants terminate (the short-route classes) the budgeted plan (ROBU) improves on the box plan (RO) in every single cell, by 4 to 10 percentage points. This confirms that pricing only a bounded number of legs at their worst case removes a large part of the box's conservatism while retaining a guarantee (Section~\ref{sec:robu}).

\green{Comment: the split-break paragraph of your version (38\,\% of oracle solutions
using at least one split break, about 0.5\,\% route-duration cost when the provision
is forbidden) is not restated here, because the \texttt{no\_split} axis has no runs
on disk: \texttt{ls solutions/} finds cs30, cs100, kw150/700/1000, kwh300/700/900 and
diesel, but nothing for no-split. The earlier flat result was in any case traced to
oracle caches carrying git conflict markers. It is a cheap run and worth redoing,
since it is the one place where the paper justifies modelling Art.~7 in full.}


\begin{table}[htbp]
\centering
\caption{Feasibility and robustness statistics by method: HoS-violation rate, SOC-violation (stranding) rate, repair frequency (offline plans), and window hit rate, computed over runs whose plan the solver certified.  ``Not opt.'' is the share of runs returned at the time limit without a proven optimum and ``Opt.\ gap'' the mean final MIP gap over those solves alone, the proven optima contributing no gap. \blue{It averages the solver logs retained, which is 178 of 178 for RO, 303 of 303 for SP, and 0 of 46 for ROBU --- hence the empty ROBU cell.}}
\label{tab:feasibility}
\resizebox{\textwidth}{!}{%
\begin{tabular}{lcccccc}
\toprule
\textbf{Method} & \textbf{HoS viol. (\%)} & \textbf{SOC viol. (\%)} & \textbf{Repairs (\%)} & \textbf{TW hit (\%)} & \textbf{Not opt. (\%)} & \textbf{Opt. gap (\%)}\\
\midrule
Greedy & \blue{8.9} & \blue{0.0} & -- & \blue{64.8} & -- & -- \\
RO & 0.0 & 0.0 & -- & \blue{50.9} & \blue{29.1} & 4.1 \\
ROBU & 0.0 & \blue{1.3} & -- & \blue{63.1} & \blue{24.2} & -- \\
LA & \blue{1.7} & \blue{4.2} & -- & \blue{76.3} & -- & -- \\
SP & 0.2 & \blue{11.8} & \blue{96.9} & \blue{84.5} & \blue{46.3} & \blue{12.2} \\
\bottomrule
\end{tabular}%
}
\end{table}


Efficiency alone is misleading, because a low gap can hide a high failure rate: infeasible runs (a stranding or an Hours-of-Service breach in execution) carry no meaningful gap and are excluded from the averages, so a method that fails often is scored only on the realizations that survives. The infeasibility strip of Figure~\ref{fig:resultsgap} and Table~\ref{tab:feasibility} restore this axis.

The two-stage plan attains the best window compliance in the study (TW hit \blue{84.5\%}) and among the lowest gaps, but only through near-permanent online recourse: it repairs its committed structure on \blue{96.9\%} of runs and still strands the vehicle on \blue{11.8\%} (by an order of magnitude the highest stranding rate of the five methods). Furthermore, \blue{46.3\%} of SP solves returned at the time limit without a proven optimum, at a mean final gap of \blue{12.2\%}, therefore providing potential low quality solution to the simulation.
The robust plans never breach the Hours-of-Service rules and never strands either, by construction, at the cost of the worst window compliance (\blue{50.9\%}) and the largest gaps. Its \blue{29.1\%} not optimal share is driven by the long-route instances, where the offline solve reaches the time limit.
ROBU trades a \blue{1.3\%} stranding rate for a \blue{12-point} gain in window hit rate (\blue{63.1\%}). However solving it in larger than short instances become intractable.
The greedy heuristic \blue{never strands (0.0\%)} but breaches driving limits on \blue{8.9\%} of runs, and misses the most windows after RO (TW hit \blue{64.8\%}), consistent with a policy that stops only when forced.
The look-ahead policy is the only method that is simultaneously highly feasible (\blue{1.7\%} HoS, \blue{4.2\%} stranding) and window-compliant (\blue{76.3\%}), which is what makes it the most defensible operational recommendation even though SP dominates it on the short routes.
\green{Comment: three things moved enough to change what this paragraph argues.
(i) Greedy's HoS-breach rate went from 4.2 to 8.9\,\% and its stranding rate to
exactly zero. That is the pre-rest-spread-cap fix in the greedy policy: it now
charges to full and stops only when genuinely forced, so it never runs the battery
down but it does overrun the 4.5\,h and shift-spread limits twice as often as
before. Greedy is now the \emph{least} HoS-compliant method in the table, which
strengthens rather than weakens your ``stops only when forced'' reading --- but the
sentence should say so explicitly instead of leading with stranding.
(ii) SP still strands an order of magnitude more than anyone else, but LA at 4.2\,\%
is no longer negligible next to ROBU's 1.3\,\%. ``Simultaneously highly feasible''
is fair on HoS and shaky on SOC; consider ``the only method that is both
HoS-compliant and window-compliant''.
(iii) ROBU's not-optimal share nearly doubled against the previous regeneration,
13.3 to 24.2\,\%, because the new medium-route ROBU attempts are counted and almost
all of them hit the limit. Its other three columns are unchanged, so the sentence
``solving it in larger than short instances become intractable'' is now supported
by this column too, not only by the runtime table.}

\begin{table}[htbp]
\centering
\caption{\red{Computation times by method and instance class: offline solve time (s) and per-stop online decision time (s).}}
\label{tab:runtime}
\begin{tabular}{lcccccc}
\toprule
& \multicolumn{3}{c}{\textbf{Offline solve (s)}} & \multicolumn{3}{c}{\textbf{Per-stop decision (s)}}\\
\cmidrule(lr){2-4}\cmidrule(lr){5-7}
\textbf{Method} & Small & Medium & Large & Small & Medium & Large\\
\midrule
Greedy & -- & -- & -- & 0.00 & 0.00 & 0.00 \\
RO & \blue{14.5} & \blue{2196.4} & \blue{7146.6} & 0.00 & 0.00 & 0.00 \\
ROBU & 2062.0 & \blue{14369.1} & -- & 0.00 & \blue{0.00} & -- \\
LA (25S-24h) & -- & -- & -- & \blue{10.76} & \blue{18.56} & \blue{20.81} \\
SP (25S) & \blue{238.6} & \blue{6145.1} & \blue{7207.0} & 0.37 & \blue{0.74} & \blue{0.52} \\
Oracle & -- & -- & -- & -- & -- & -- \\
\bottomrule
\end{tabular}
\end{table}
% NOTE: the Oracle offline solve time is not stored in the oracle_*.json cache, so those cells show '--'.


Finally, the methods differ by orders of magnitude in cost, and along two distinct dimensions: an offline solve committed before departure and an online decision cost paid at every stop (Table~\ref{tab:runtime}). Greedy is effectively instantaneous ($<10^{-4}$\,s per stop). The offline plans pay their cost upfront and then execute for free. While SP and RO manage to solve short and medium route instances within the timelimit, they are struggling to solve completely on the long route instances. \blue{ROBU now terminates on the medium routes as well, but at 14\,369\,s of wall clock, roughly seven times the short-route figure and twice the per-solve time limit, and it still returns nothing on the long routes.}
LA carries the opposite profile: no offline solve, but \blue{11--21}\,s of scenario optimization \emph{per stop}, growing only mildly with route length. That is comfortably real-time against the 60\,km charger spacing of the base case (a decision every $\approx$40\,min of driving), \blue{and it stays cheaper in total than the offline plans: on a long route with $\approx$170 stops LA accumulates about 3\,500\,s, roughly half of what RO spends offline.}
\green{Comment: two corrections here. (i) ``ROBU cannot solve any instance larger than
the small route instances'' is contradicted by tex/tables/runtime.tex, which now
carries a medium-route ROBU offline time of 14\,369.1\,s, and by
tex/tables/gap.tex, which carries a single medium/Few/None ROBU cell at 52.3\,\%
over $n=1$ (24 further medium/Few/None attempts and all 15 medium/Few/Tight ones
came back unsolved). I left that cell as ``--'' in Table~\ref{tab:gap} since one run
is not a class average, but the runtime claim had to be softened. Note that 14\,369\,s
exceeds the 7\,200\,s limit stated in \S\,8.2, because ROBU's scenario generation is a
sequence of master solves and the limit applies per solve --- worth stating.
(ii) The closing sentence claimed LA spends more total wall clock than RO. It does
not, and did not in the previous draft either: $170 \times 20.8 \approx 3\,500$\,s against
RO's 7\,146\,s, and the measured LA wall-clock median on the long routes
(data\_output/additional\_la\_stats.csv) is 3\,444\,s. I reversed it.}


\subsection{Additional analyses}
\label{sec:additional-protocol}

The base-case instances of Section~\ref{sec:basecase} are too expensive to re-run for every sensitivity, so the analyses of Sections~\ref{sec:sensitivity}--\ref{sec:diesel} use a reduced set. All additional experiments run on the short, medium and long route classes with few and many customers, the None and Tight window classes, and seeds 1--25, giving 300 instances.
The additional analysis use Greedy, LA and the oracle while SP, RO and ROBU are omitted.
\blue{The pack-against-charger grid of Section~\ref{sec:grid} is the one exception: it
runs the Greedy policy alone, on the short and medium route classes only, over the
nine crossed cells that the one-at-a-time sweeps do not already cover, for a
further 1\,800 simulated routes.}


\subsubsection{Look-ahead sensitivity}
The rolling-horizon look-ahead method is parameterized through several options. First, the look-ahead horizon, which determines how far ahead we look. While the base case considers a horizon of 24 hours, additional testing is done with both shorter and longer horizons. Second, the number of scenarios considered when selecting an action provides information about the distribution of uncertainty and greater insight into potential extreme scenarios. Finally, to achieve a reasonable solve time compatible with real-life use, the base case considers only the LP relaxation of the subproblem, whereas here we examine solving the full MILP subproblems.

\begin{figure} [h!]
    \centering
    \includegraphics[width=1\linewidth]{figures/resultsLA.png}
    \caption{Look-ahead configuration analysis, one axis per column: horizon at the base scenario count (left) and scenario count at the base horizon (right). The strip below each column shades the infeasibility rate per route class}
    \label{fig:resultsLA}
\end{figure}




{\color{red}
We observe in Figure \ref{fig:resultsLA} that while increasing the number of scenarios or the horizon length increases significantly the decision time per stop to to either more subproblems or bigger subproblems, the gap to oracle does not decrease except when going from 12h horizon to 24h horizon. Indeed, allowing each subproblem to see 24h in advance instead of 12h permit to anticipate the next place of rest, considering the cycle of duty that is roughly 24h. However, seing 2 days in advance does not procudre more advantages. For the number of scenarios, while increasing it does not give better results in term of gap, it secure the solution chosen in term of infeasibility rate. Seeing more scenarios allows to see potential extreme cases that necessiate immediate actions and discard risky actions, providing a safer schedule.}
\blue{Concretely, halving the horizon to 12\,h costs 1.2--1.7 percentage points of
median gap on every route class (short 4.3 to 5.5 without windows and 7.6 to 9.7
with tight ones) while doubling it to 48\,h returns nothing (4.4 and 8.2), at
three times the decision time on the long routes (44--50\,s against 20--21\,s).}
\green{Comment: the second half of the paragraph, on scenario count, is not supported
by the current runs. Ten scenarios give the same median gap as twenty-five
(short 4.2/7.7 against 4.3/7.6, medium 4.3/6.1 against 4.6/6.8, long 3.8/5.0
against 3.9/5.3) \emph{and} the same infeasibility: 2 infeasible runs out of 100 on
the short routes at $|\Xi|=10$ against 0 at 25, 1 out of 100 on the medium against
4, 7 out of 97 on the long against 5 of 100. Fifty scenarios are no safer either
(1 of 100 short, 2 of 100 medium, 1 of 46 long) and cost 27--42\,s per stop. What
\emph{does} drive
infeasibility is the horizon, not the scenario count: the 12\,h horizon fails on
9\,\% of short, 18\,\% of medium and 21\,\% of long runs against 0/4/5\,\% at 24\,h,
which is the same duty-cycle argument you already make for the gap. I would move
the safety claim onto the horizon axis and drop it from the scenario axis.}

{\color{red}
We also analyse the runs with MILP subproblems. As expected, the results improve as
the gap to oracle reduces from \blue{4.3--7.6\%} to \blue{2.1--2.9\%} on the short routes,
and from \blue{4.6--6.8\%} to \blue{2.3--4.1\%} on the medium ones. However rthe decision time
per stop explodes and reach \blue{48--61\,s on the short routes and 84--119\,s on the medium
ones, an order of magnitude above the LP relaxation}, which is unreallistic for a real
life use by a truck driver. }
\green{Comment: the previous draft quoted ``6.7\,\% to 4.7\,\%'' and ``almost 500\,s per
stop''; neither matches tex/tables/additional\_la\_policy.tex. The MIP tail is now
much cheaper than that (a median of 48--119\,s) but also much more effective
(it roughly halves the gap rather than shaving a fifth off it), so the trade-off
you describe is real but the arithmetic behind ``unrealistic'' is weaker: at
$\approx$40\,min of driving between stops, 120\,s is not obviously out of reach. The
argument that survives cleanly is the total wall clock, 8\,300--12\,000\,s per medium
route. Also note MIPTAIL still has no long-route runs at all, and that the medium
tight cell moved from 3.6 to 4.1\,\% between the two regenerations while its decision
time went from 103 to 119\,s --- the medium MIPTAIL runs grew from 34 to 48 and the
new ones are the harder end, so treat that cell as still settling.}


\subsubsection{Sensitivity analysis on charging infrastructure}
\label{sec:sensitivity}

The base case fixes a 500\,kWh battery charged on 350\,kW charger, and charging stations every 60\,km as mandated for the TEN-T core network by Regulation~(EU) 2023/1804.  We perturb these three parameters one at a time.

\begin{figure}[htbp]
    \centering
    \includegraphics[width=\linewidth]{figures/add_sens.png}
    \caption{One-at-a-time sensitivity of mean route duration to the charging
    infrastructure, relative to the base case, over the paired instances}
    \label{fig:sens}
\end{figure}





Halving the CS spacing from 60 to 30\,km buys only 1.1\% of route duration at the hindsight optimum, while doubling it to 100\,km, which is the network density advised for the secondary road network, costs 1.5\%: at the mandated spacing the corridor is already dense enough that the binding constraint is how long a charging stop takes, not how far away the next one is.
Upgrading the charger from the 350\,kW base to 700\,kW saves \blue{4.2\%}, and going all the way to a 1\,MW megawatt-charging system saves \blue{5.3\%}, so a 700\,kW charger already captures the majority of the benefit available. Downgrading to 150\,kW costs a significant 18.6\% of route duration, highlighting the need for dedicated heavy duty truck chargers. Indeed, at 150 kW the charge no longer fits inside the mandatory break, and every charging event starts adding time to the schedule instead of hiding inside a break the driver has to take anyway.
Investment should therefore go into more power per charger in each CS rather than into a denser than necessary charging infrastructure.
Then, shrinking the pack from 500 to 300 kWh costs 6.3 \% at the optimum and \blue{15.0} \% under the greedy policy, while growing it to 700 or 900 kWh returns only \blue{2.4} and \blue{3.3} \%.
\green{Comment: the two incomplete rows improved but are not fixed. The 900\,kWh oracle
figure is now 3.3\,\% over $n=199$ (medium $-3.85$ and long $-2.73$, still no short
column), up from 2.7\,\% over $n=79$ long-route-only in the previous regeneration ---
so that number moved by 0.6 points purely from coverage, and will move again when
the short routes land. The 1\,MW figure is still long-route-only, $n=84$ of a nominal
300. Every other row has $n=300$ across the three classes. The 1\,MW cell is the one
I would prioritise, because the conclusion leans on it for the ``megawatt chargers
buy little'' claim, and because 5.3 against 4.2\,\% at 700\,kW is a 1.1-point
increment that a route-class reweighting could easily halve.}

Two secondary observations are worth noting.
First, the effect is larger on the medium and large routes than on the short ones on every power level, because a longer route accumulates more charging events over which the saving adds up, whereas the spacing effect is almost route-length independent.
Second, \blue{Greedy and LA bracket the oracle on the power downgrade rather than
under-reacting to it uniformly: Greedy loses 19.7\,\% and LA 15.2\,\% against the
oracle's 18.6\,\%.}
\green{Comment: this second observation needs rewriting, not just renumbering. In the
version you sent both policies under-reacted (14.5 and 15.3 against 18.6). After the
greedy fix, Greedy now \emph{over}-reacts to the 150\,kW downgrade (19.7) and, more
importantly, no longer ``tracks the oracle closely on the battery axis'' at all:
at 300\,kWh Greedy loses 15.0\,\% and LA 11.6\,\% against the oracle's 6.3\,\%, i.e.\ more
than twice the optimal cost. The mechanism is the one you already describe
elsewhere --- a myopic charge-to-full policy on a small pack stops more often and
strands the schedule against the break structure --- and it is a better story than
the old one, since it says the value of a big pack is mostly \emph{organisational}.
LA also over-reacts badly on the 100\,km spacing axis (+6.3\,\% against the oracle's
+1.5\,\%). Suggested replacement: ``the two online policies exaggerate every
degradation, by a factor of about two on the battery axis and four on the spacing
axis, so an operator without a planning tool sees a harsher infrastructure
dependence than the corridor physically imposes''.}


\subsubsection{Battery capacity against charger power}
\label{sec:grid}

The two axes of Section~\ref{sec:sensitivity} that matter, pack size and charge point
power, cannot be read independently. The minimum state of charge $E^{\min}$ and the
point at which the charging curve tapers both scale with $E^{\mathrm{cap}}$, so moving
the pack alone moves range and taper avoidance together, and a one-at-a-time sweep
cannot say which of the two produced the duration change. We therefore cross the two
axes on a $4 \times 4$ grid, from 300 to 900\,kWh against 150 to 1\,000\,kW. Reading
across a row isolates charge speed at a fixed pack, reading down a column isolates
pack at a fixed charge point, and any curvature between the two is the interaction
that the one-at-a-time sweeps fold away. The base row (350\,kW) and the base column
(500\,kWh) are exactly the sweeps of Figure~\ref{fig:sens} and share their instances
and runs with them, so only the nine crossed cells are new.

Two restrictions keep the grid affordable. First, it is run under the Greedy policy
alone. Read on route duration rather than on the penalised objective, Greedy sits
within about two percentage points of LA in every cell of Table~\ref{tab:gap}, and it
reproduces the sign and the ordering of every axis in Figure~\ref{fig:sens} except
the megawatt cell discussed at the end of this section, while
costing no solver time at all --- which matters when nine cells at 200 instances
each is already 1\,800 simulated routes. Second, it is run on the short and medium
route classes only. The long class behaves as an extrapolation of the other two on
every axis rather than as a separate regime --- the spacing effect is within one
percentage point of the short and medium values, and the power effect is monotone in
route length (14.4, 19.4 and 21.9\,\% at 150\,kW) with no reordering --- while carrying
by far the slowest solves. It keeps its one-at-a-time sweeps in Figure~\ref{fig:sens}.

\begin{figure}[htbp]
    \centering
    % produced by `python -m src.plot.additional_figures --section grid`:
    % figures/additional_grid_battery_power.pdf ; per-cell n in
    % data_output/additional_grid_stats.csv
    \includegraphics[width=\linewidth]{figures/grid_battery_power.png}
    \caption{Battery capacity against charger power, short and medium routes.
    Each cell is the mean paired change in Greedy route duration against the base
    case (500\,kWh, 350\,kW, boxed), over the instances that are feasible on both
    sides of the pair.}
    \label{fig:grid}
\end{figure}

Figure~\ref{fig:grid} shows that the two levers are not substitutes. The 150\,kW
column is costly at every pack size and stays costly: 14.9\,\% at the base pack on the
short routes and 14.8\,\% at 900\,kWh, 20.4 and 20.6\,\% respectively on the medium ones.
Tripling the battery recovers essentially nothing from an under-powered charge point,
which is the grid's version of the asymmetry already visible in
Section~\ref{sec:sensitivity}. The 300\,kWh row is costly at every power level in the
same way: 13.6\,\% at the base charge point on the short routes, and still 11.9\,\% at
1\,MW, so a megawatt charger recovers under two percentage points of what the small
pack costs. A range shortfall cannot be bought back with power any more than a power
shortfall can be bought back with range, because the small pack forces additional
stops whose access and manoeuvre time no charge rate removes.

Where the two axes do interact is in the interior, and mildly in favour of doing both.
Going from 500 to 900\,kWh at the base charge point buys 2.7\,\% on the short routes and
2.3\,\% on the medium ones; going from 350 to 700\,kW at the base pack buys 2.6 and
3.3\,\%; doing both buys 5.8 and 6.3\,\%, marginally more than the sum of the two. The
practical reading is that pack and charge point should be specified together, and that
the cheapest configuration reaching a target duration is not the largest value of
either axis on its own.

One cell breaks monotonicity: at the base 500\,kWh pack, the 1\,MW charge point is
\emph{worse} than the 700\,kW one (+0.5 and +0.4\,\% against $-2.6$ and $-3.3$\,\%). The same
sign appears in the one-at-a-time sweep, where Greedy reports $+0.3$\,\% at 1\,MW while
the oracle reports $-5.3$\,\%, so this is a property of the myopic policy rather than of
the corridor: once the charge is already shorter than the break it runs alongside,
extra power buys no makespan, and a charge-to-full rule spends the surplus on the
tapered tail of the curve, which the higher nominal power does not accelerate.
\green{Comment on the grid, three things. (i) The premise that Greedy is a good proxy
holds for \emph{ranking} cells but not for their magnitudes: on the battery axis Greedy
loses 15.0\,\% at 300\,kWh where the oracle loses 6.3\,\%, so every number in
Figure~\ref{fig:grid} should be read as what a myopic operation would measure, not as
the optimal cost of that configuration. Worth one sentence in the text, and worth
re-running the four corner cells under the oracle if you want the figure to carry a
policy claim. (ii) The 1\,MW / 500\,kWh anomaly is reproducible across both route
classes and both the grid and the sweep, so it is not seed noise, but I have not
verified the taper explanation against the charging curve --- please check
TAIL\_C\_RATE before the sentence goes in. A related oddity: on the medium routes the
300\,kWh row is slightly \emph{worse} at 700\,kW (+14.9) than at 350\,kW (+14.3), which the
taper story does not explain. (iii) Per-cell $n$ ranges from 73 to 98 rather than the
nominal 100, because the pairing drops any instance that was infeasible on either
side, and the 300\,kWh rows strand most --- so the small-pack cells rest on the fewest
pairs and are the least comparable. The counts are in
data\_output/additional\_grid\_stats.csv if you want them as a footnote.}


\subsubsection{Comparison with diesel trucks}
\label{sec:diesel}

Every instance of the subset in Section~\ref{sec:additional-protocol} is re-solved as a diesel vehicle. Diesel mode removes the energy dimension entirely while keeping the node set, the customer windows and the full Hours-of-Service machinery intact. Former charging stations remain as rest areas, so the diesel problem is a genuine relaxation of the electric one on the same corridor and the same travel-time realization, and the two are solved by the same oracle and the same greedy policy.

\begin{figure}[htbp]
    \centering
    % \red{PLACEHOLDER} - produced by `python additional_figures.py --section diesel`:
    % figures/additional_diesel_gap.pdf
    \includegraphics[width=0.75\linewidth]{figures/diesel.png}
    \caption{Electrification penalty in route duration, per route class:
    the penalty actually realized by the greedy policy and by the hindsight
    optimum.  The annotation below each group gives the mean coupling share
    $\sum_i g_i/\sum_i \tau^c_i$ and the number of instances.}
    \label{fig:diesel}
\end{figure}





Electrification carries a real but moderate makespan penalty at the hindsight optimum: $+7.8\%$ on the short routes, $+8.9\%$ on the medium and $+9.8\%$ on the long ones. Break/charge coupling works as the model predicts, with 70--76\% of charging time credited inside a mandatory break, but it does not cancel the cost of charging.

Decomposing where the penalty goes, total charging time is the whole of the gross cost (+2.7, +6.4 and +11.9 h), and the largest single credit against it is break time no longer taken separately (-1.3, -2.9 and -4.4 h): the coupling effect of \citet{brockmann_break-and-charge_2025}, recovered here at roughly \blue{37 to 49} \% of the gross charging cost. What that work does not model is the rest of the table: queuing, the net stop maneuver and the repositioning off the charging bay together add +0.7, +1.3 and +2.5 h, which is roughly half of the coupling credit.
\green{Comment: ``roughly 37\,\%'' is the long-route value only; the ratio of the coupling
credit to gross charging time is 49\,\% on the short routes, 45\,\% on the medium and
37\,\% on the long, so it falls with route length rather than sitting at one number.
Everything else in this subsection is unchanged against the new runs --- the
decomposition table, the coupling shares and the three headline penalties all
reproduce exactly.}

\green{Comment: the electrification penalty is also the one place where the long-route
oracle certificate matters. tex/tables/additional\_diesel.tex reports a worst-case
remaining MIP gap on the EV side of 0.5\,\% (short), 3.4\,\% (medium) and 9.2\,\% (long)
--- a column that is not currently in the paper and probably should be --- and since
the diesel side closes to tolerance, the long-route $+9.8\,\%$ is an \emph{upper} bound on
the true optimal penalty. Worth one sentence, because it means the 8--10\,\% headline
in the abstract is conservative rather than optimistic.}

\section{Conclusion} \label{section:conclusion}

% restate problem
This paper studied the joint scheduling of driver activities and partial non-linear fast charging for a battery electric truck on a fixed long-haul route under travel-time uncertainty. We formulated a mixed-integer linear program capturing the daily provisions of EU Hours-of-Service regulation and validated its regulatory core against an independent exact method. Around this model we built five decision architectures spanning the range from full offline commitment to full online re-optimization, evaluated them in a common discrete-event simulator over \blue{900} instances, and benchmarked them against a perfect-hindsight oracle.
\green{Comment: 950 does not correspond to anything on disk. The benchmark is the
$3 \times 3 \times 4 \times 25 = 900$ instances of Table~\ref{tab:instances}, plus the one real-life
corridor of Section~\ref{sec:casestudy}; coverage per method is uneven within that
(Table~\ref{tab:gap}), which is a different statement and belongs in \S\,8.2 rather
than here. Note also that the abstract says ``six decision architectures'' counting
the oracle, and this sentence says five plus the oracle --- pick one.}

% findings
{\color{red}
The main findings are as follows. First, the methods that adapt en route outperform those that commit a full schedule at departure, across every instance class. When evaluated in the simulation framework, the committed plans yield both longer route durations and lower time-window compliance, and the second effect is the larger of the two: a schedule fixed before departure absorbs travel-time uncertainty in its rest structure but fails in the compliance to time windows.}

{\color{red}
Second, part of the duration cost of electrification comes not from charging itself but from access to the charging station, maneuvering and queuing, and this component is amplified by the number of stops a longer route requires. Mandatory breaks do supply usable charging windows, and our schedules exploit \blue{70--76\%} of them, but the residual charging and access time still lengthens the route relative to a diesel vehicle on the same corridor. Interestingly, a myopically scheduled operation observes \blue{roughly 40\,\% more than} this penalty, so about \blue{a third} of the electrification cost an operator would measure is attributable to scheduling rather than to charging.}
\green{Comment: ``roughly twice'' overstates it. Greedy against the hindsight optimum is
11.6 vs 7.8, 12.8 vs 8.9 and 13.5 vs 9.8\,\% by route class, i.e.\ a factor of
1.4--1.5, so the schedulable share is around a third, not a half. The qualitative
claim --- that a meaningful slice of what an operator would call an electrification
cost is organisational --- is unchanged and is, if anything, easier to defend at a
third.}

Finally, a sensitivity analysis reveals that investing in more powerful chargers reduces route duration more than increasing the density of the charging network. The binding constraint on these corridors is how long a charging stop takes, not how far away the next one is: charging duration is what binds, and additional charger power allows the battery to be fully recharged within the mandatory break. At the other end of the range, megawatt chargers deliver little beyond intermediate charger sizes: once the charge is shorter than the break it runs alongside, further power buys nothing at all, because the break must be taken regardless. The asymmetry matters here a lot: a 150 kW corridor costs 19 \% of route duration where a corridor thinned to 100 km spacing costs 1.5 \%, so an under-powered network cannot be rescued by building more of it. \blue{Crossing charger power against battery capacity shows the converse to be equally true: a 900\,kWh pack recovers essentially nothing from a 150\,kW corridor, and a megawatt charge point recovers under two percentage points from a 300\,kWh pack, so the two must be specified together rather than traded against each other.}

These findings can be set against the charging-operator perspective of \citet{pagnier_charging_2026}, where maximizing the route coverage of a charging network favors a few very large stations over many smaller ones. As outlined there, the problem then becomes one of grid access, on infrastructure that is already frequently saturated.

% work direction
Several directions could reduce the limitations of the present work. Correlating travel times between legs, and letting the level of uncertainty differ between near and distant legs, may yield different results: current GPS tracking allows congestion to be estimated fairly accurately over a short horizon, whether it arises from peak-hour effects or from recurring patterns.
Next, endogenous charger availability would close the gap between our deterministic queueing time and the behaviour a shared corridor actually produces.
Furthermore, a cost objective incorporating electricity pricing would change the trade-off between charging early and charging cheaply, and could yield a different solution structure as prices vary by location and time of day. Since driver wages remain the dominant cost component, however, we expect a cost-minimising schedule to resemble the duration-minimising one studied here.
Finally, relaxing the fixed customer sequence would let the vehicle trade a routing detour against a better-placed charging opportunity, which is the degree of freedom this study deliberately removes.



\appendix
\section{Constraints linearization} \label{section:appendix1}
The constraints relative to linearization of constraints in the deterministic MILP formulation are summarized here.

\begin{alignat}{2}
    l_i^1 &\leq T^{\mathrm{drv}}_{\mathrm{cons}}\, r_i,
    &\quad& \, i \in \mathcal{I} \label{eq:l1_ub1}\\
    l_i^1 &\leq c_i^d,
    &\quad& \, i \in \mathcal{I} \label{eq:l1_ub2}\\
    l_i^1 &\geq c_i^d - T^{\mathrm{drv}}_{\mathrm{cons}}\bigl(1-r_i\bigr),
    &\quad& \, i \in \mathcal{I} \label{eq:l1_lb}
\end{alignat}



\begin{alignat}{2}
    l_i^2 &\leq T^{\mathrm{drv}}_{\mathrm{sh_2}}\, \rho_i,
    &\quad& \, i \in \mathcal{I}, \label{eq:l2_ub1}\\
    l_i^2 &\leq s_i^d,
    &\quad& \, i \in \mathcal{I}, \label{eq:l2_ub2}\\
    l_i^2 &\geq s_i^d - T^{\mathrm{drv}}_{\mathrm{sh_2}}\bigl(1-\rho_i\bigr),
    &\quad& \, i \in \mathcal{I}. \label{eq:l2_lb}
\end{alignat}




\begin{alignat}{2}
    l_i^4 &\leq T^{spr}_{2}\, \rho_i,
    &\quad& \, i \in \mathcal{I}, \label{eq:l4_ub1}\\
    l_i^4 &\leq s_i^w,
    &\quad& \, i \in \mathcal{I}, \label{eq:l4_ub2}\\
    l_i^4 &\geq s_i^w - T^{spr}_{2}\bigl(1-\rho_i\bigr),
    &\quad& \, i \in \mathcal{I}. \label{eq:l4_lb}
\end{alignat}


\begin{alignat}{2}
    l^5_i &\leq T^{spr}_2\, \rho_i, &\qquad& i \in \mathcal I \label{eq:l5_1}\\
    l^5_i &\leq h_i + o_i, &\qquad& i \in \mathcal I \label{eq:l5_2}\\
    l^5_i &\geq h_i + o_i - T^{spr}_2\,(1 - \rho_i), &\qquad& i \in \mathcal I \label{eq:l5_3}
\end{alignat}


\section{Valid inequalities for the Rolling-horizon subproblems} \label{sec:valid_inequ}

Constraints \eqref{eq:vi-2} fix the upper bound on the number of charging stop on each subsection of the path to what is required due to battery constraints. Constraints \eqref{eq:vi-3}--\eqref{eq:vi-4} ensure that each subsection of the route has the minimum required number of break or rest based on driving time.

\begin{equation}
    e^d_i-e^a_i \leq (E^{cap}-E^{min})y_i,
    \qquad \, i \in \mathcal{K}.
    \label{eq:vi-1}
\end{equation}
\begin{equation}
    \sum_{l \in \mathcal K, l<i} y_l \geq \lceil \frac{1}{E^{cap}-E^{min}}(\sum_{l < i} E_l - (E^{0}-E^{min})) \rceil,
    \qquad \, i \in \mathcal{K}.
    \label{eq:vi-2}
\end{equation}
\begin{equation}
    \sum_{l<i} \rho_l \geq \lceil \frac{\sum_{l<i} D_l}{T^{drv}_{sh2}} \rceil-1,
    \qquad \, i \in \mathcal{I}.
    \label{eq:vi-3}
\end{equation}
\begin{equation}
    \sum_{l<i} r_l \geq \lceil \frac{\sum_{l<i} D_l}{T^{drv}_{cons}} \rceil -1,
    \qquad \, i \in \mathcal{I}.
    \label{eq:vi-4}
\end{equation}

\begin{comment}


\section*{CRediT authorship contribution statement}
\noindent
\textbf{Céline Pagnier:} Conceptualization, Methodology, Software, Formal Analysis, Visualization, Writing - original draft.
\textbf{Steffen J. Bakker:} Conceptualization, Methodology, Software, Formal Analysis, Visualization, Writing - original draft.
\textbf{Pernille Seljom:} Conceptualization, Methodology, Software, Writing - Review \& Editing.
\textbf{Asgeir Tomasgard:} Conceptualization, Methodology.


\section*{Acknowledgements}
\noindent
This paper was prepared as a part of the Norwegian Research Centre for Energy
Transition Studies (FME NTRANS, p-nr: 296205), funded by the Research Council of Norway.
\end{comment}


\bibliography{paper4}

\end{document}

\endinput
%%
%% End of file `elsarticle-template-num.tex'.
